% !TEX root = ../nc-cubic.tex
\subsubsection{Hesse 境界の線束・残余対象と低次積の積分模型}
\label{subsubsec:nfm-hesse-boundary-models}

\begin{lemma}[追加境界の八直線と一つの付値]
\label{lem:nfm-period-hesse-eight-boundary-lines}
$\zeta$ を原始三乗根、$\epsilon=\zeta^2-\zeta$ とし、
\[
 D=1-\frac{s}{24}-\frac{t}{216}-\frac{s^2}{6912}
\]
と置き、Hesse 不変量写像を $r(p,q)$ と書く。その引戻しは
\[
 D\bigl(r(p,q)\bigr)
 =-\frac1{6912}(q^2-4)
       \prod_{i=0}^{2}\bigl((2\zeta^ip-q)^2+12\bigr).
\]
従って $D=0$ の逆像は、相異なる八本の直線
\[
 q=\pm2,\qquad 2\zeta^ip-q=\pm2\epsilon\quad(i=0,1,2)
\]
からなる。
$q=2$ の一般点で Hesse 写像は非分岐であり、$v_D$ の延長の一つは
$v_{q-2}$ で分岐指数一である。
$q=2+z$ と書くと
\[
 D\bigl(r(p,2+z)\bigr)
 =-\frac{(p^3+8)^2}{27}z+O(z^2).
\]
したがって、既存の二境界仮定の下で
$\overline\gamma_n=[D]^{b_n}$ と分かっているなら、
$b_n$ は Hesse 周期 torsor の $q=2$ における mod-three 付値に等しい。
\end{lemma}
\begin{proof}
積の式は $\zeta^2+\zeta+1=0$ を用いた多項式恒等式である。
$\epsilon^2=-3$ により表示した一次因子へ分解する。
既存の Jacobian の式を $q=2$ に制限すると
\[
 \det Dr=-384p^2(p^3+8)^2\ne0,
 \qquad \Delta=-64p^3(p^3+8)^3\ne0
\]
である。これにより一般点での非分岐性が従う。
また、その線上では
\[
 s=16-16p^3,\qquad t=64+160p^3-8p^6.
\]
剰余体は $\mathbb C(p)/\mathbb C(p^3)$ という三次拡大であり、
この三次性は分岐指数ではない。
$D$ の一次係数を計算すると主張の式となる。
分岐指数が一なので、$K_0$ 上の $D$-付値と制限後の $q-2$-付値は一致する。
\end{proof}

周期位相の値とは独立に、この境界の点線束と三分点被覆は次のように記述できる。

\begin{proposition}[Hesse 境界の成分群と三分点の局所分解]
\label{prop:nfm-hesse-boundary-third-division}
$k_0=\mathbb C(p,m)$ とし、$p$ と任意個数の混合・純中心係数 $m$ を
独立変数とする。$R=k_0[[z]]$、$K=k_0((z))$、$q=2+z$ と置き、
\[
 A=p(1-q^2/4),\qquad B=q+p^3/3+q^3/12,
 \qquad
 \mathcal E:\ \frac A3(X^3+Y^3+Z^3)+BXYZ=0
\]
を考える。$E=\mathcal E_K$ の原点と点 shift を
\[
 O=(1:-1:0),\qquad P=(1-q/2:-1-q/2:-p)
\]
と取る。このとき次が成り立つ。
\begin{enumerate}
\item $\mathcal E/R$ は正則な split $I_3$ 模型であり、
特殊 fibre は $C_X\cup C_Y\cup C_Z=V(XYZ)$ である。
$O$ は $C_Z$ の滑らかな点、$P$ は $C_X$ の滑らかな点へ特殊化する。
従って $P$ の N\'eron 成分群 $\mathbb Z/3$ における像は非零である。
\item $\mathcal M=\mathcal O_{\mathcal E}(P-O)$ は実際の可逆層で、
特殊 fibre の multidegree は $(1,0,-1)$ である。
同じ一般 fibre の線束を延ばして multidegree $(0,0,0)$ にすることはできない。
一方、$\mathcal M^3\otimes\mathcal O_{\mathcal E}(C_X-C_Z)$ の
multidegree は $(0,0,0)$ である。
\item $\zeta$ を原始三乗根とし、
\[
 T=(0:1:-1),\qquad U=(1:-\zeta:0)
\]
と置く。これらは定数群 $E[3]$ の基底であり、
Weil pairing による Kummer 類の二座標は
\begin{equation}
 \label{eq:nfm-hesse-boundary-kummer-values}
 \bigl(e_{3,*}(\kappa_3(P),T),e_{3,*}(\kappa_3(P),U)\bigr)
   =\left(\left[\frac{p^3+8}{2z}\right],1\right)
   \quad\text{in }(K^\times/K^{\times3})^2.
\end{equation}
従って
\begin{equation}
 \label{eq:nfm-hesse-boundary-division-fields}
 [3]^{-1}(P)\simeq\coprod_{j=1}^{3}\operatorname{Spec}L,
 \qquad
 L=K\left(\sqrt[3]{\frac{2z}{p^3+8}}\right).
\end{equation}
$L/K$ は完全分岐の巡回三次拡大である。tame inertia は
$\langle U\rangle$ の非零平行移動として作用し、九点上の cycle type は $3^3$ である。
慣性軌道の集合は split な階数三の非分岐商となる。
任意の一つの三分点を定義する有限拡大の分岐指数は三の倍数である。
\end{enumerate}
この幾何的慣性型は、一般の active 基底と混合係数を加えた
標準 split 係数空間の $D=0$ の一般点でも同じである。
これは周期 torsor $T_{\eta,n}$ の分岐を計算する主張ではない。
\end{proposition}
\begin{proof}
補題\ref{lem:nfm-generic-hesse-third-division}の逆座標により、表示した $P$ は
本稿の右点対応の平行移動点である。直接代入でも $P\in E(K)$ を確認できる。
\[
 A=-pz+O(z^2),\qquad B=(p^3+8)/3+O(z)
\]
であり、$p(p^3+8)$ は単元である。三角形の各頂点では、方程式の
$z$ 微分が $-p/3$ となるので全空間は正則である。
その他の特殊 fibre の点は相対的にも滑らかである。
三成分はそれぞれ二つの他成分と横断的に交わり、その自己交点数は $-2$ だから、
これは極小の split $I_3$ 模型である。
その滑らかな部分が N\'eron 模型を与え、成分群は $\mathbb Z/3$ である
\cite[Appendix~C, Theorems~15.1--15.2]{NFMSilvermanArithmetic}。
$P_0=(0:-2:-p)$ と $O_0=(1:-1:0)$ は異なる成分の滑らかな点である。

この二切断は相対 Cartier 因子なので $\mathcal M$ は可逆であり、
multidegree は表示の通りである。
同じ一般 fibre を持つ二つの可逆層の比は、正則模型上の垂直因子
$a_XC_X+a_YC_Y+a_ZC_Z$ で与えられる。
その multidegree の変化は
\[
 (-2a_X+a_Y+a_Z,\ a_X-2a_Y+a_Z,\ a_X+a_Y-2a_Z)
\]
である。$(1,0,-1)$ を零にするには、第一成分と第二成分の差から
$3(a_X-a_Y)=1$ が必要となり不可能である。
三乗後には $C_X-C_Z$ の multidegree $(-3,0,3)$ がこれを相殺する。

局所分解を定数も含めて計算する。
Hesse 曲線の九つの flex はすべて $K$-有理であり、$O$ を原点に取ると
それらが $E[3]$ をなす。
$T$ は非零成分、$U$ は単位成分の非零三捩れ点なので独立である。
三つの接線を
\[
 \ell_O=A(X+Y)-BZ,\qquad
 \ell_T=-BX+A(Y+Z),\qquad
 \ell_U=AX+A\zeta^2Y-B\zeta Z
\]
と書く。$O$ の局所座標 $w=Z/X$ に対して
\[
 X=1,\qquad
 Y=-1+\frac BA w-\frac{A^3+B^3}{3A^3}w^3+O(w^4)
\]
だから、$\ell_O=-(A^3+B^3)w^3/(3A^2)+O(w^4)$ である。
従って $\ell_T/\ell_O$、$\ell_U/\ell_O$ の $w^{-3}$ 係数はそれぞれ
\[
 c_T=\frac{3A^2(A+B)}{A^3+B^3},\qquad
 c_U=-\frac{3A^3(1-\zeta^2)}{A^3+B^3}
\]
であり、$F_T=(\ell_T/\ell_O)/c_T$、$F_U=(\ell_U/\ell_O)/c_U$ と置く。

この先頭係数一の規格化が Kummer 計算に必要である。
例えば $\operatorname{div}(F_T)=3(T)-3(O)$ であり、
$[3]^*((T)-(O))$ は $K$-有理な主因子なので、それを因子に持つ
$G_T\in K(E)^\times$ を取れる。
$G_T=a w^{-1}+O(1)$ とすると、$w\circ[3]=3w+O(w^2)$ より
\[
 F_T\circ[3]=(3a)^{-3}G_T^3.
\]
従って $F_T\circ[3]$ は $K(E)$ における三乗であり、
Tate--Lichtenbaum pairing の評価公式
\cite[Chapter~XI, Proposition~9.2]{NFMSilvermanArithmetic}から
$e_{3,*}(\kappa_3(P),T)=[F_T(P)]$ を得る。$U$ についても同様である。

接線の表示へ $P$ を代入して整理すると
\begin{align*}
 F_T(P)&=
 \left(\frac{p+2}{3p}\right)^3
 \frac{p^3+8}{2z}\bigl(1+O(z)\bigr),\\
 F_U(P)&=
 \frac{\zeta}{3(1-\zeta^2)}
 \left(\frac{p^3+8}{3p}\right)^3z^{-3}
 \bigl(1+O(z)\bigr).
\end{align*}
$1+zR$ の元は一意に $1+zR$ 内で三乗根を持ち、
第二式の残る定数は $\mathbb C^\times$ の元なので三乗である。
これが式\eqref{eq:nfm-hesse-boundary-kummer-values}を与える。

$E[3]$ は定数群であり、Weil pairing は非退化である。
従ってこの torsor の Galois 像は
$\ker e_3(-,U)=\langle U\rangle$ 全体で、その固定体は
式\eqref{eq:nfm-hesse-boundary-division-fields}の $L$ である。
$v_z(2z/(p^3+8))=1$ だから $L/K$ は完全分岐の巡回三次拡大となる。
九点への作用は位数三の平行移動なので、その軌道は三つである。
各点の定義体は $L$ であり、分解と分岐指数の主張が従う。
慣性軌道集合は $E[3]/\langle U\rangle$ の定数 torsor である。

最後に一般係数空間への移行を確認する。
$q=2$ の一般点では $g=p\sum X^3/3+qXYZ$ 自身は滑らかであり、
$\operatorname{SL}_3$ 軌道の接空間は八次元である。
補題\ref{lem:nfm-period-hesse-eight-boundary-lines}の $\det Dr\ne0$ により、
Hesse の二変数方向は不変量の二つの接方向を張る。
従って $\operatorname{SL}_3\times\mathbb A^2_{p,q}$ から全十係数空間への
写像はこの一般点で \'etale である。
独立な混合・純中心係数を加えても同じであり、必要な frame の追加は滑らかな底変更である。
この非分岐な chart で幾何的 inertia が検出されるから、
一般係数空間でも同じ $3^3$ の慣性型を持つ。
ここでは一般三分点被覆の代数的な点曲線公式だけを用い、
full 残余比較の境界への延長を仮定していない。
\end{proof}

\begin{lemma}[境界の点線束の線型行列表示]
\label{lem:nfm-hesse-boundary-determinantal-line}
前命題の特殊 fibre を $C_0=V(XYZ)$、$\mathcal M_0=\mathcal M|_{C_0}$ と書く。
その cohomology は
\[
 (h^0,h^1)(\mathcal M_0(j))=
 \begin{cases}
 (3j,0)&j\ge1,\\
 (0,0)&j=0,\\
 (0,-3j)&j\le-1.
 \end{cases}
\]
$i_0:C_0\hookrightarrow\mathbb P^2_{k_0}$ に対し、完全列
\begin{equation}
 \label{eq:nfm-hesse-boundary-determinantal-line}
 0\longrightarrow\mathcal O(-1)^3
 \xrightarrow{\ L_0\ }\mathcal O^3
 \longrightarrow i_{0,*}\mathcal M_0(1)\longrightarrow0,
 \qquad
 L_0=\begin{pmatrix}Z&0&0\\pY&-2Z&X\\-X&Y&0\end{pmatrix}
\end{equation}
が存在し、$\det L_0=-XYZ$ である。
さらに、この完全列は $\mathcal E\hookrightarrow\mathbb P^2_R$ 上の
$\mathcal M(1)$ の相対的な三行三列の線型行列表示へ持ち上がる。
\end{lemma}
\begin{proof}
各成分上の次数は $(j+1,j,j-1)$ である。
$j=0$ では、次数 $-1$ の成分上の切断が零であり、節点での貼合せから
次数零の成分上でも零となる。残る次数一の成分では二つの節点で零となるため、
切断は零である。Euler 標数が零なので $H^1$ も零である。
$j\ge1$ では双対線束の次数が $(-j-1,-j,-j+1)$ となり、
同じ節点での議論によりその $H^0$ が零である。
$\omega_{C_0}\simeq\mathcal O_{C_0}$ と Serre 双対性により
$H^1(\mathcal M_0(j))=0$ となり、Riemann--Roch から $h^0=3j$ を得る。
$j\le-1$ も同様である。

$\mathcal M_0(1)$ の三つの切断を、各成分上の有理切断として
\[
 \begin{array}{c|ccc}
 C_X&-pY^2/(2Z-pY)&YZ/(2Z-pY)&2Z^2/(2Z-pY)\\
 C_Y&X&0&Z\\
 C_Z&X+Y&0&0
 \end{array}
\]
で与える。分母は $P_0$ で許された一位の極を表し、
$X+Y$ は $O_0$ で必要な零点を持つ。
三節点で値が一致するのでこれらは大域切断であり、各成分上で基点を持たない。
従って $\mathcal O^3\to i_{0,*}\mathcal M_0(1)$ は全射である。
$L_0$ の各列がこれらの切断の関係式であることは成分ごとの代入で分かる。
行列式は $-XYZ$ であり、二次小行列式に $X^2$、$-2Z^2$、$pY^2-2XZ$ が含まれるので、
$L_0$ は $C_0$ の全点で階数二である。
従って cokernel は $C_0$ 上の可逆層であり、上の全射は
$\operatorname{coker}L_0\simeq i_{0,*}\mathcal M_0(1)$ を与える。

相対表示については、proper cohomology と底変更により
$\pi:\mathcal E\to\operatorname{Spec}R$ に対する
$\pi_*\mathcal M(1)$ が自由階数三で、その基底を上の三切断から持ち上げられる。
その評価射の cokernel は特殊 fibre 上で零である。
支持の proper 性と Nakayama の補題により評価射は全体でも全射である。
$i_*\mathcal M(1)$ は Cartier 因子上の可逆層の直像なので局所射影次元一を持ち、
この評価射の核 $\mathcal K$ は階数三のベクトル束である。
底変更により $\mathcal K|_{\mathbb P^2_{k_0}}\simeq\mathcal O(-1)^3$ となる。
再び cohomology と底変更を $\mathcal K(1)$ へ適用すると、
その三つの基底切断を持ち上げられ、評価射
$\mathcal O^3\to\mathcal K(1)$ は特殊 fibre 上の同型を持ち上げる。
退化 locus の proper 性からこれは全体で同型である。
従って相対的にも $\mathcal K\simeq\mathcal O(-1)^3$ となり、
特殊 fibre が $L_0$ である線型行列表示を得る。
\end{proof}

\begin{lemma}[境界の点線束の有理式による相対表示]
\label{lem:nfm-hesse-boundary-explicit-line}
命題\ref{prop:nfm-hesse-boundary-third-division}の記法を用い、
\[
 F_E=\frac A3(X^3+Y^3+Z^3)+BXYZ,\qquad
 b=\frac{p^3+8}{3},\qquad c=\frac{p(p^3+q^3)}3
\]
と置く。以下の式は完備化前の DVR
$R_0=k_0[q]_{(q-2)}$ 上で定義される。
\begin{align*}
 x_f&=X+Y,&y_f&=X-Y,&w_f&=A(X+Y)-BZ,\\
 u&=-q/c,&v&=2/c,\\
 a_0&=-\frac{4A(A^3+B^3)}{3B^3},&
 a_2&=1+\frac{4A^3}{B^3},&
 a_4&=-\frac{4A^2}{B^3},\qquad
 a_6=\frac{4A}{3B^3}.
\end{align*}
行列
\[
 N=\begin{pmatrix}
 x_f-uw_f&y_f-vw_f&0\\
 0&a_0x_f+(a_0u+a_2)w_f&w_f\\
 y_f+vw_f&(a_0u^2+a_2u+a_4)w_f&-x_f
 \end{pmatrix},
\]
および定数行列
\[
 U=\begin{pmatrix}4&-2&-2p\\-4bp&0&0\\4&-2&2p\end{pmatrix},
 \qquad
 V=\begin{pmatrix}0&-p&-1\\2p&p&-1\\2p&-p&1\end{pmatrix}
\]
に対して $L=U^{-1}NV$ と置く。このとき
\[
 L\bmod(q-2)=L_0,\qquad
 \det L=-F_E/b,
\]
であり、完全列
\[
 0\longrightarrow\mathcal O_{\mathbb P^2_{R_0}}(-1)^3
 \xrightarrow{\ L\ }\mathcal O_{\mathbb P^2_{R_0}}^3
 \longrightarrow i_*\mathcal O_{\mathcal E}(P-O)(1)
 \longrightarrow0
\]
を得る。従って、完備化や三分点の添加を必要としない有限な有理式で、
前補題の特定の特殊 fibre 表示を持ち上げられる。
また
\[
 Q=-b\operatorname{adj}(L),\qquad LQ=QL=F_E I_3
\]
は指定した方程式 $F_E$ に対する相対 graded matrix factorization である。
これは曲線上の点線束の模型であり、full 比較核やその周期位相の計算を
それだけで与えるものではない。
\end{lemma}
\begin{proof}
$B,c,b$ は $q=2$ でそれぞれ $b,pb,b$ へ移り、すべて単元である。
表示した座標変換の行列式は $2B$ なので $R_0$ 上で可逆である。
直接展開により
\[
 4F_E=y_f^2w_f-a_0x_f^3-a_2x_f^2w_f-a_4x_fw_f^2-a_6w_f^3,
 \qquad v^2=a_0u^3+a_2u^2+a_4u+a_6
\]
を得る。$P$ の新座標は $(-q:2:c)=(u:v:1)$、
$O=(0:1:0)$ である。後者を用いた三行三列の行列式の展開は
$\det N=4F_E$ を与える。
また
\[
 \det U=-32bp^2,\qquad\det V=8p^2,\qquad
 N|_{q=2}V=UL_0
\]
であるから、行列の特殊化と determinant の式が従う。

一般 fibre の線束を符号も含めて確認する。
$\bar x=x_f/w_f$, $\bar y=y_f/w_f$ と書くと、
$N^{\mathsf t}$ の左核は
\[
 \left(\frac{\bar y-v}{\bar x-u},\ -1,\
              a_0(\bar x+u)+a_2\right)
\]
で与えられる。第一成分は $-P=(u:-v:1)$ で一位の極、
$O$ で高々一位の極を持ち、第三成分は $O$ で二位の極を持つ。
三成分は独立であり、次数三の可逆層
$\mathcal O_E((-P)+2O)$ の全切断空間の基底となる。
次数三の線束は genus-one 曲線上で大域生成されるので、
これらの切断と行列の関係式から
\[
 \operatorname{coker}N^{\mathsf t}
   \simeq i_*\mathcal O_E((-P)+2O)
   \simeq i_*\mathcal M^{-1}(1)
\]
を得る。ここで $\mathcal O_E(1)=\mathcal O_E(3O)$ と
$(-P)+P\sim2O$ を使った。
線型表示を $\mathcal O(-1)$ で双対化すると、三次 Cartier 因子に対する
\[
 \mathcal E xt^1(i_*\mathcal N,\mathcal O(-1))
       \simeq i_*(\mathcal N^\vee(2))
\]
により $\operatorname{coker}N\simeq i_*\mathcal M(1)$ となる。
この転置が点 $P$ の符号を決める。

相対的には、$L_0$ の階数が $C_0$ の全点で二である。
階数が一以下となる閉集合は射影空間上で proper であり、
特殊 fibre と交わらないので空である。従って $\operatorname{coker}L$ は
$\mathcal E$ 上の可逆層 $\mathcal N$ の直像である。
一般 fibre 上で $\mathcal N\simeq\mathcal M(1)$ だから、
正則模型上で両者の比は垂直因子
$a_XC_X+a_YC_Y+a_ZC_Z$ による。
特殊 fibre でも $L=L_0$ により両者は同じ multidegree $(2,1,0)$ を持つ。
$I_3$ の交点行列の核は $(1,1,1)$ のみなので、
$a_X=a_Y=a_Z$ である。全 fibre は主因子であり、
$\mathcal N\simeq\mathcal M(1)$ が相対的に従う。
最後の積の式は adjugate 恒等式である。
\end{proof}

\begin{proposition}[周期を変えない残余三次の積分模型]
\label{prop:nfm-hesse-boundary-integral-residual-model}
命題\ref{prop:nfm-hesse-boundary-third-division}の
$R=k_0[[z]]$, $K=k_0((z))$, $r=n-1$ を用いる。
良い一般係数体からこの Hesse 体へ移した残余三次、
指定同型 $\varphi_{\mathrm{sc}}$ および full 尺度を、
既存の固定入力の下で取る。
点三次の指定生成元を
\[
 J_R=\frac{A}{3}(X^3+Y^3+Z^3)+BXYZ
\]
と書く。このとき、基礎体を拡大せずに、次を満たす射影的平坦模型
$\mathcal Y\subset\mathbf P_R^{n+1}$ を取れる。
\begin{enumerate}
\item 一般 fibre は指定された $Y_{\mathrm{res}}/K$ と同型であり、
      active 平面切断は厳密に $\mathcal E=V(J_R)$ である。
      選んだ模型に対し、$\varphi_{\mathrm{sc}}$ は点平面と残余平面の
      $R$ 上の同型へ延びる。
\item 方程式 $F_R$ の一般 fibre は cube-frame の生成元を
      線型座標変更で移したものである。
      その変更は active 平面で恒等であり、方程式に追加の scalar を掛けない。
      従って指定 MF 周期と一般 fibre の位相類は変わらない。
\item $\mathcal C:=\mathcal Y\cap V(u_1,\ldots,u_r)=\mathcal E$ は
      $\mathcal Y$ 内の長さ $r$ の正則列で定義される。
      $\mathcal Y$ の全空間は $\mathcal C$ に沿って正則である。
      点線束 $\mathcal M=\mathcal O_{\mathcal E}(P-O)$ と
      補題\ref{lem:nfm-hesse-boundary-determinantal-line}の線型表示が
      この残余切断上へ移る。特に $i_*\mathcal M$ は
      $\mathcal Y$ 上の perfect 複体である。
\end{enumerate}
ここで模型は選択に依存する。特殊 fibre の滑らかさ、既存の good 開への所属、
元の残余 moduli 射の境界への延長、または full 比較核の延長は主張しない。
\end{proposition}

\begin{proof}
まず周囲空間と方程式を固定する。
$\varphi_{\mathrm{sc}}$ は active な射影平面の同型なので、残余の
Brauer--Severi 周囲空間は split な平面の $K$-点を含む。
従って周囲空間も split である。
この事実は、Brauer--Severi 空間の $K$-点が対応する中心単純代数の
最小右イデアルを与え、その作用表示が行列代数への同型を与えることによる。
$\varphi_{\mathrm{sc}}$ の線型持上げを全周囲空間の線型 frame へ延長し、
中心方向の補空間を選ぶ。

この frame における cube-frame の生成元を $F\in K[X,Y,Z,u]_3$ とする。
命題\ref{prop:nfm-point-normalized-jacobian-cover}の尺度同定と
式\eqref{eq:nfm-period-cube-frame}を active 平面へ制限すると、
$F|_{u=0}\otimes P_{\mathrm{src}}=J_R\otimes P_{\mathrm{src}}$ である。
従って $F|_{u=0}=J_R$ が厳密に成り立つ。
ここでは任意の生成元を後から $K$-scalar で規格化したのではなく、
点尺度が定める指定生成元を用いている。

$u$ 次数によって $F=J_R+F_1+F_2+F_3$ と分ける。
各 $F_j$ の係数は有限個なので、非負整数 $N$ を十分大きく取り、
全ての非零係数 $c$ に対して
\[
 v_z(c)+Nj\ge0
\]
とできる。そこで
\[
 a_N:(X,Y,Z,u)\longmapsto(X,Y,Z,z^N u),\qquad
 F_R=a_N^*F=J_R+\sum_{j=1}^3z^{Nj}F_j
\]
と置くと、$F_R\in R[X,Y,Z,u]_3$ である。
$a_N$ は一般 fibre 上の線型同型で、active 平面では恒等である。
$B\bmod z=(p^3+8)/3\ne0$ なので $F_R\bmod z$ は非零である。
従って $z$ は $R[X,Y,Z,u]/(F_R)$ 上非零因子であり、
この商およびその射影化は $R$-flat である。
一般 fibre の同型と active 切断の同定もこの式から従う。

座標変更 $a_N$ は曲率を厳密に $F\mapsto a_N^*F$ と移す。
graded MF の底変更は内部 shift と複体 shift をそのまま保つので、
指定された三次周期も厳密に移す。
特に $\det(a_N|_{K_{\mathrm{pt}}})=1$ であり、
cube-frame の determinant 因子も変化しない。
最初の active 線型持上げを $\lambda$ 倍した場合の差は
命題\ref{prop:nfm-period-equation-cube-frame}の三乗倍だけであり、
位相類とその mod-three 付値を変えない。
この操作に $z$ の根の添加はない。

任意の順序で $u_1,\ldots,u_r$ を零にしても、残る方程式は
非零な項 $J_R$ を含む。従って各段階で次の $u_i$ は
方程式を割らず、多項式環の一意分解性から商の非零因子である。
これで正則列性を得る。同じ議論は閉 fibre と一般 fibre にも通る。
その閉埋込み $i$ は有限 Tor 次元を持ち、$\mathcal M$ は
$\mathcal C$ 上可逆だから $i_*\mathcal M$ は perfect である。
全空間の正則性は、特殊 fibre の三頂点以外では active 偏微分が非零であり、
三頂点では $\partial_zF_R|_{u=0}=\partial_zJ_R=-p/3$ が単元であることによる。
一般 fibre は元の滑らかな残余三次である。
平面と曲線の同定により、線束とその行列表示も移る。
\end{proof}

\begin{lemma}[点線束の直像の相対 Ext は平坦である]
\label{lem:nfm-hesse-boundary-pushforward-ext-flat}
命題\ref{prop:nfm-hesse-boundary-integral-residual-model}の模型において、
$i:\mathcal C\hookrightarrow\mathcal Y$、$\pi_C:\mathcal C\to\operatorname{Spec}R$、
$\pi_Y:\mathcal Y\to\operatorname{Spec}R$ と書く。
法束を $N_{\mathcal C/\mathcal Y}\simeq\mathcal O_{\mathcal C}(1)\otimes W$、
$W=R^r$ と自明化する。
任意の整数 $a,b,j$ と $d=b-a$ に対し、自然な完全列
\begin{align*}
 0&\longrightarrow
 R^1\pi_{C,*}\mathcal O_{\mathcal C}(d+j-1)
       \otimes\bigwedge^{j-1}W\\
 &\longrightarrow
 H^jR\pi_{Y,*}R\mathcal Hom_{\mathcal Y}
       (i_*\mathcal M(a),i_*\mathcal M(b))\\
 &\longrightarrow
 \pi_{C,*}\mathcal O_{\mathcal C}(d+j)
       \otimes\bigwedge^jW\longrightarrow0
\end{align*}
がある。ただし $q\notin\{0,\ldots,r\}$ なら $\bigwedge^qW=0$ とする。
これらの相対 Ext は有限自由であり、任意の $R$-底変更に両立する。
とくに $a=b$ の階数は
\[
 \operatorname{rank}\operatorname{Ext}^j_{\mathcal Y/R}(i_*\mathcal M,i_*\mathcal M)
 =\begin{cases}
 1&j=0,\\
 1+3r&j=1,\\
 3j\binom rj&2\le j\le r,\\
 0&\text{その他}
 \end{cases}
\]
である。完全列の正準分裂や dg 代数の形式性は主張しない。
\end{lemma}

\begin{proof}
$i$ は正則閉埋込みで $\mathcal M$ は可逆だから、局所 Koszul 分解により
\[
 \mathcal Ext^q_{\mathcal Y}
       (i_*\mathcal M(a),i_*\mathcal M(b))
 \simeq i_*\!\left(
       \mathcal O_{\mathcal C}(d+q)\otimes\bigwedge^qW\right).
\]
これは線束 $\mathcal M$ の局所自明化に依存しない。
相対曲線の cohomological dimension は一なので、相対 local-to-global
Ext spectral sequence の非零列は $p=0,1$ だけである
\cite[Tag~0BQP]{NFMStacksRelativeExt}。
従って $E_2$ 以降の微分は全て零となり、表示した完全列を得る。

任意の $k$ に対し、$\mathcal O_{\mathcal C}(k)$ の閉 fibre と一般 fibre の
cohomology の次元は共に
\[
 (h^0,h^1)=
 \begin{cases}
 (3k,0)&k>0,\\
 (1,1)&k=0,\\
 (0,-3k)&k<0
 \end{cases}
\]
である。滑らかな三次では genus 一の Riemann--Roch による。
三角形では負次数の切断が各成分上で零であり、
$H^0(\mathcal O)=k_0$、$\omega_{C_0}\simeq\mathcal O_{C_0}$ なので、
Serre 双対性と Riemann--Roch から同じ式を得る。
proper cohomology と底変更により、表示した二つの外側の項は有限自由で、
その形成は底変更に両立する。従って中央の相対 Ext も有限自由である。

$\mathcal Y/R$ は proper な平坦 hypersurface で、二つの直像は perfect である。
従ってその相対 derived Hom は perfect であり、derived 底変更に両立する
\cite[Lemma~36.30.1]{NFMStacksRelativeExt}。
これは有限 affine cover の \v{C}ech 複体で計算しても確認できる。
得られた cohomology が全て $R$-free なので、任意の底変更について
cohomology の Tor 項は消える。これで通常の相対 Ext の底変更も従う。
最後の階数は $d=0$ を完全列へ代入したものである。
\end{proof}

\begin{remark}[残余射影を施した低 Ext との区別]
上の平坦性は $i_*\mathcal M(a)$ 同士の相対 Ext についての主張である。
$q_Y^R\mathcal M=j_Y^!i_*\mathcal M$ とその rotation 軌道へ移るには、
例外対象による射影の cone と、その中の合成写像の階数を別に調べる必要がある。
free な cohomology を持つ有限複体の間の写像でも、cone に torsion が生じうる。
従ってこの補題だけから周期評価に必要な $S_p$ の低 Ext 平坦性や
三次関係の原始性を結論しない。
\end{remark}

\begin{lemma}[三角形切断に沿う滑らかさの有限判定]
\label{lem:nfm-hesse-boundary-transverse-smoothness}
$k$ を標数零の体、$c\in k^\times$ とし、
\[
 F_0=cXYZ+\sum_{i=1}^r u_iQ_i(X,Y,Z)+O((u)^2),
 \qquad r=n-1
\]
を三次式とする。$\mathcal Y_0=V(F_0)$ が
$C_0=V(u_1,\ldots,u_r,XYZ)$ に沿って幾何学的に滑らかであるための
必要十分条件は、三頂点 $e_X,e_Y,e_Z$ のそれぞれについて
\[
 (Q_1(e_\alpha),\ldots,Q_r(e_\alpha))\ne0
 \qquad(\alpha=X,Y,Z)
\]
となることである。
また、active 制限を $cXYZ$ に固定した全三次式の affine 空間には、
全体が幾何学的に滑らかな三次超曲面を定める非空開がある。
\end{lemma}

\begin{proof}
$C_0$ の三頂点以外では $XYZ$ の active 偏微分の少なくとも一つが非零である。
頂点では active 偏微分が全て零で、残りの偏微分はちょうど
$\partial_{u_i}F_0=Q_i(e_\alpha)$ である。Jacobian 判定から最初の主張を得る。

非空性は代数閉包上で確認すればよい。
固定平面を $\Lambda$ とし、$W=H^0(\mathbf P^{n+1},I_\Lambda(3))$ とする。
$a\notin\Lambda$ では、ある $u_i(a)\ne0$ であり、
$u_i H^0(\mathcal O(2))\subset W$ は $a$ の全一次 jet を生成する。
従って $F=cXYZ+w$ が $a$ で特異となる条件は、$w\in W$ への
$n+2$ 個の独立な affine 線型条件である。
$a$ を $n+1$ 次元の周囲空間の $\Lambda$ 外で動かしても、
対応する incidence の次元は $\dim W-1$ 以下である。
従ってその像の閉包は $W$ の真の閉集合である。
一方、三頂点での上の非零条件もそれぞれ非空な開条件であり、
例えば $u_1(X^2+Y^2+Z^2)$ の係数を用いて全てを同時に満たせる。
この有限個の非空開の共通部分を取れば所要の非空開を得る。
\end{proof}

\begin{remark}[積分模型と滑らかな残余模型の差]
命題\ref{prop:nfm-hesse-boundary-integral-residual-model}で
全ての不等式を $v_z(c)+Nj>0$ とすると、特殊方程式は
$F_0=((p^3+8)/3)XYZ$ になる。
これは三頂点で上の横断ベクトルが全て零となる特異な模型である。
従って pole を消すだけでは、滑らかな残余三次の good reduction を得ない。
他方、三角形切断自体は滑らかな全三次の模型を妨げないことも前補題が示す。
残る具体的な問題は、実際の残余係数に適合する積分 lattice を選び、
三頂点の横断ベクトルと切断外の Jacobian を検査することである。
一般混合係数が存在するというだけで、この特定の模型の閉 fibre が
上の非空開に入るとは結論できない。
\end{remark}

\begin{proposition}[点線束から残余対象への相対右射影]
\label{prop:nfm-boundary-right-projection}
命題\ref{prop:nfm-hesse-boundary-integral-residual-model}の記法で
$r=n-1$、$\pi:\mathcal Y\to\operatorname{Spec}R$、
$i:\mathcal C\hookrightarrow\mathcal Y$ とする。
$\mathcal M$ は移送した点線束である。
次の構成は、$\mathcal Y$ の特殊 fibre の滑らかさを仮定しない。
相対 perfect 圏内で
\[
 \mathcal A_R=\langle\mathcal O_{\mathcal Y},\ldots,
              \mathcal O_{\mathcal Y}(r-1)\rangle^\perp
\]
と置き、その包含を $j_R$ と書く。
この包含は右随伴 $j_R^!$ を持ち、
\begin{equation}\label{eq:nfm-boundary-right-projection}
 j_Rj_R^!F=
 R_{\mathcal O(-1)}R_{\mathcal O(-2)}\cdots
 R_{\mathcal O(-r)}F
 \qquad(F\in\operatorname{Perf}\mathcal Y)
\end{equation}
である。右辺は右端から適用し、相対 right mutation の規約は
\[
 R_EF\longrightarrow F\longrightarrow
 \pi^*\bigl(R\pi_*R\mathcal Hom(F,E)\bigr)^\vee\otimes E
 \longrightarrow R_EF[1]
\]
とする。
従って
\[
 \mathcal Z_R=j_R^!i_*\mathcal M[1]
\]
は、一般 fibre で周期評価の対象 $Z_p=q_Y^RM[1]$ に一致する
実際の相対 perfect 対象を与える。
この構成と、その有限個の twist、mutation、Hom の導来複体は、
閉 fibre と一般 fibre への導来底変更に両立する。

最初の right mutation $R_{\mathcal O(-r)}i_*\mathcal M$ は自明である。
特に $n=3$ の場合には次の一つの coevaluation だけでよい：
\begin{equation}\label{eq:nfm-boundary-threefold-one-cone}
 j_Rj_R^!i_*\mathcal M\longrightarrow i_*\mathcal M
 \longrightarrow\mathcal O_{\mathcal Y}(-1)\otimes_R V^\vee[2],
 \qquad
 V=\pi_{\mathcal C,*}\mathcal M^{-1}(1)\simeq R^3.
\end{equation}
この相対対象の構成は、比較核、低 Ext の平坦性、
または source と target の周期の一致を主張しない。
\end{proposition}

\begin{proof}
$\mathcal Y$ は平坦な Cartier 三次超曲面なので、
proper な完全射であり、相対双対化複体は
$\omega_{\mathcal Y/R}[n]=\mathcal O_{\mathcal Y}(-r)[n]$ である。
有限表示の射影埋込みと通常の超曲面完全列から
\[
 R\pi_*\mathcal O_{\mathcal Y}=R,\qquad
 R\pi_*\mathcal O_{\mathcal Y}(-d)=0\quad(1\le d\le r-1)
\]
を得る。この計算は両 fibre に同じである。
従って表示した正の line bundle block と
$\langle\mathcal O(-r),\ldots,\mathcal O(-1)\rangle$ は
相対 exceptional block である。
各評価・余評価の cone は perfect 対象を perfect 対象へ移す。

相対双対性により、perfect 対象 $G$ について
\[
 R\pi_*R\mathcal Hom(G,\mathcal O(-j))
 \simeq
 \bigl(R\pi_*R\mathcal Hom(\mathcal O(r-j),G)\bigr)^\vee[-n]
 \qquad(1\le j\le r).
\]
よって $\mathcal A_R$ は負の block の左直交でもある。
負の block を $\mathcal O(-r),\ldots,\mathcal O(-1)$ の順で
右変異すると、この左直交への射影を得る。
既に消した Hom は、line bundle block の半直交性により後続の変異でも零を保つ。
射影の triangle に $\mathcal A_R$ の対象から Hom を施せば、
この射影が包含の右随伴であることが分かる。
これで式\eqref{eq:nfm-boundary-right-projection}を得る。

$i$ は $r$ 本の中心一次式による regular immersion であり、
$\mathcal M$ は可逆なので $i_*\mathcal M$ は $\mathcal Y$ 上 perfect である。
従って上の有限変異を実際に適用できる。
一般 fibre で $\mathcal A_K=\operatorname{Ku}(Y)$ であり、
右随伴の一意性によって $j_R^!i_*\mathcal M$ の一般 fibre は
$j_Y^!i_*M=q_Y^RM$ となる。
ここで $\mathbf Li^*$ の右随伴は $i_*$ であり、
$i^!$ や余法線の twist を $q_Y^R$ 自体へ挿入しない。

具体的な余評価空間を求めるときには regular immersion の双対性を用い、
\begin{equation}\label{eq:nfm-boundary-coevaluation-space}
 R\pi_*R\mathcal Hom_{\mathcal Y}(i_*\mathcal M,\mathcal O(-j))
 =R\pi_{\mathcal C,*}\bigl(\mathcal M^{-1}(r-j)\bigr)[-r]
\end{equation}
を得る。ここで normal determinant は $\mathcal O_{\mathcal C}(r)$ である。
境界の multidegree 計算を $\mathcal M^{-1}$ にも適用すると、
$R\pi_{\mathcal C,*}\mathcal M^{-1}=0$、
$R\pi_{\mathcal C,*}\mathcal M^{-1}(1)$ は自由階数三である。
特殊 fibre 上の消滅から、properness と Nakayama により同じ相対主張が従う。
従って $j=r$ の余評価空間は零である。
$r=2$ では残る $j=1$ の空間が $V[-2]$ なので、
式\eqref{eq:nfm-boundary-threefold-one-cone}を得る。

proper perfect 射による perfect 複体の押し下げは $R$-perfect であり、
ここで使った pullback、tensor、双対、評価の finite cone は導来底変更に両立する。
これにより各 Hom 複体も $R$ 上の有界有限自由複体で表せる。
その cohomology が自由であることは、perfect 性からは従わない。
\end{proof}

\begin{proposition}[三次元境界の低 orbit Hom の平坦性]
\label{prop:nfm-boundary-threefold-low-hom}
命題\ref{prop:nfm-boundary-right-projection}で $n=3$ とし、
$E=i_*\mathcal M$、$Q=j_Rj_R^!E$、$Z=Q[1]$ と置く。
$\rho=L_{\mathcal O}\circ(1)$ と $\sigma=\rho^{-1}$ を
相対残余 perfect 圏の rotation とする。
次の四つの Hom 複体の cohomology は $R$-自由で、両 fibre への底変更と可換である：
\[
\begin{array}{c|c}
 q& H^iR\operatorname{Hom}(Z,\rho^qZ)\text{ の非零階数}\\\hline
 0&(i,\operatorname{rk})=(0,1),(1,10)\\
 -1&(i,\operatorname{rk})=(1,5),(2,5)\\
 -2&(i,\operatorname{rk})=(2,10),(3,1)\\
 -3&(i,\operatorname{rk})=(2,1),(3,10).
\end{array}
\]
さらに、実際の integral 方程式 $F$ が定める周期 class は
$H^2R\operatorname{Hom}(Z,\rho^{-3}Z)\simeq R$ の生成元である。
この結果は $m_2,m_3$ の値や、それらから回復した三次 relation の原始性を
まだ決定しない。
\end{proposition}

\begin{proof}
まず射影前の Hom を計算する。
$i$ は二本の中心一次式による regular immersion で、その normal bundle は
$N_{\mathcal C/\mathcal Y}=\mathcal O_{\mathcal C}(1)^2$ である。
可逆層の自己 Ext で $\mathcal M$ は相殺するので
\[
 \mathcal Ext^j_{\mathcal Y}(E,E(d))
   =i_*\bigl(\mathcal O_{\mathcal C}(d+j)\otimes\Lambda^jR^2\bigr).
\]
曲線の cohomological dimension は一である。
従って local-to-global spectral sequence は退化し、
各 $a$ について
\[
 0\to H^1(\mathcal C,\mathcal O(d+a-1))\otimes\Lambda^{a-1}R^2
 \to\operatorname{Ext}^a(E,E(d))
 \to H^0(\mathcal C,\mathcal O(d+a))\otimes\Lambda^aR^2\to0
\]
を得る。$\mathcal C$ は平面三次の平坦族であり、
全ての $\mathcal O_{\mathcal C}(k)$ の cohomology は、
平面の超曲面完全列から自由で両 fibre に同じである。
従って上の列の全項は自由で底変更と可換となる。
特に $d=0$ の Ext 階数は次数 $0,1,2$ で $(1,7,6)$、
$d=-1$ の Ext 階数は次数 $1,2$ で $(5,5)$ である。

$A=\mathcal O_{\mathcal Y}(-1)$、
$V=H^0(\mathcal C,\mathcal M^{-1}(1))$、
$W=H^0(\mathcal C,\mathcal M(1))$ と置く。
$V,W$ は自由階数三である。
式\eqref{eq:nfm-boundary-threefold-one-cone}と右随伴の性質により
\[
 R\operatorname{Hom}(Q,Q)
 =R\operatorname{Hom}(Q,E)
 =\operatorname{Cone}\bigl(
     (V\otimes W)[-2]\xrightarrow{\mu}R\operatorname{Hom}(E,E)\bigr).
\]
regular immersion の双対性で最高 local Ext を同定すると、
$H^2(\mu)$ は
\[
 V\otimes W\longrightarrow H^0(\mathcal C,\mathcal O(2))
\]
という線束の乗法写像である。

この写像の特殊 fibre を、既存の点線束行列 $L_0$ で検査できる。
$L_0^t$ の cokernel は $\mathcal M_0^{-1}(1)$ の表示を与え、
二つの cokernel の pairing は $\operatorname{adj}L_0$ で与えられる。
実際、$L_0\operatorname{adj}L_0=\operatorname{adj}L_0L_0=-XYZ I$ なので
この pairing は cokernel へ降り、曲線上で $L_0$ の階数が二であることから
線束の完全な pairing になる。
直接計算で
\[
 \operatorname{adj}L_0=
 \begin{pmatrix}
 -XY&0&0\\-X^2&0&-XZ\\pY^2-2XZ&-YZ&-2Z^2
 \end{pmatrix}.
\]
$p$ と $2$ は単元なので、その成分は六つの二次単項式を張る。
従って $\mu$ は特殊 fibre で全射であり、Nakayama により $R$ 上でも全射となる。
その核は自由階数三である。
上の cone の長完全列から $\operatorname{Ext}^0(Q,Q)=R$、
$\operatorname{Ext}^1(Q,Q)\simeq R^{10}$、他の Ext は零となり、
底変更との両立も従う。

次に逆 rotation を用いる。
相対窓の内部 shift に対応する関手は
\[
 \sigma=R_{\mathcal O(-1)}\circ(-1).
\]
この式と $Q=R_{\mathcal O(-1)}E$ から
\[
 \sigma Q=R_{\mathcal O(-1)}R_{\mathcal O(-2)}E(-1)=j_Rj_R^!E(-1)
\]
を得る。
しかも $R\operatorname{Hom}(A,E(-1))=R\Gamma(\mathcal C,\mathcal M)=0$ なので、
右射影の triangle を使えば
\[
 R\operatorname{Hom}(Q,\sigma Q)
 =R\operatorname{Hom}(Q,E(-1))
 =R\operatorname{Hom}(E,E(-1)).
\]
先の $d=-1$ の計算により $q=-1$ の行が従う。

相対 Serre 関手は $S_{\mathcal Y}=(-2)[3]$ である。
左射影を $p_R$ と書くと
\[
 S_{\mathcal A_R}^{-1}=p_RS_{\mathcal Y}^{-1}j_R
 =\rho^2[-3],\qquad S_{\mathcal A_R}=\sigma^2[3].
\]
これは二つの line bundle を除く有限 mutation と相対双対性の式であり、
特殊 fibre の滑らかさを用いない。
従って
\[
 R\operatorname{Hom}(Q,\sigma^2Q)
 =R\operatorname{Hom}(Q,Q)^\vee[-3],
\]
から $q=-2$ の行が従う。

最後に $q=-3$ と指定周期を確認する。
$S/(F)$ は次数零部分が regular DVR $R$、Gorenstein parameter が二の
noetherian graded Gorenstein 環である。
Burke--Stevenson\cite[Theorem~6.4(1), $i=1$]{NFMBurkeStevenson}の窓は、
滑らかさを仮定せず
\[
 D^b(\operatorname{coh}\mathcal Y)
 =\langle\mathcal O(-2),\mathcal O(-1),\mathcal B_1\rangle
\]
と graded singularity category の同値を与える。
窓の隣接する切断を、次数一つの自由 summand の除去 triangle で比較すると、
内部 shift $(-1)$ は上の $R_{\mathcal O(-1)}\circ(-1)$ に移る。
この比較は section の truncation と実際の evaluation/coevaluation で構成される。
有限 mutation は perfect 対象を保つので、相対 perfect 残余圏にも制限できる。
中央商 $S/(F)$ の実際の自己交叉二拡張、または $D^2=+F$ の MF 表示により
$(-3)\simeq[-2]$ である。
同じ $F$ を用いる窓比較により
$\sigma^3\simeq[-2]$ とその指定周期が $R$ 上で得られる。
従って $q=-3$ の Hom 複体は $R\operatorname{Hom}(Q,Q)[-2]$ であり、
周期 class は $\operatorname{id}_Q$ に対応する。
$\operatorname{End}(Q)=R$ なのでこの class は primitive な生成元である。
これで全ての行を得る。共通の $[1]$ は二つの Hom の引数で相殺する。
\end{proof}

\begin{lemma}[点線束の線型表示からの明示的な三次 MF]
\label{lem:nfm-boundary-line-mf}
上の integral 方程式を $F\in S=R[X,Y,Z,u_1,\ldots,u_r]_3$、
$H=F|_{u=0}$ とする。
点線束の相対線型表示を $L$ とし、
$\det L=cH$、$c\in R^\times$ と書く。
$Q=c^{-1}\operatorname{adj}L$ と置けば $LQ=QL=H I_3$ である。
$F-H=\sum_{a=1}^r u_aG_a$、$G_a\in S_2$ と分解する。
このとき
\[
 (L,Q)\mathbin{\widehat\otimes}
       \widehat{\bigotimes}_{a=1}^r(u_a,G_a)
\]
は、各 parity の階数が $3\cdot2^r$ の integral graded MF of $F$ を与える。
これは $i_*\mathcal M(1)$ の graded module の stabilization であり、
source 非可換代数との Morita 比較ではない。
また、この raw MF を無指定の window で $q_Y^RM[1]$ と同一視しない。
後者の対象は命題\ref{prop:nfm-boundary-right-projection}の右射影と最後の $[1]$ で固定する。
\end{lemma}

\begin{proof}
$N_0=\operatorname{coker}(L:R[X,Y,Z](-1)^3\to R[X,Y,Z]^3)$ とする。
$H$ が正則なので $L$ は単射である。
$u_a$ を零として作用させた $S$-加群 $N$ は、
\[
 P_\bullet=[S(-1)^3\xrightarrow{L}S^3]
                  \otimes_S K_\bullet(u_1,\ldots,u_r)
\]
を有限自由分解に持つ。従って
$P_i=S(-i)^{3\binom{r+1}{i}}$ である。
その微分を $d$ と書く。
二項分解上で $Q$ を逆向きに作用させる写像を $h_H$ とし、
Koszul 外積作用を $\varepsilon_a$ と書く。
tensor の符号規約で
\[
 h=h_H+\sum_aG_a\varepsilon_a
\]
と置けば、各項は homological degree $+1$、internal degree $3$ を持ち、
\[
 d^2=h^2=0,\qquad dh+hd=F\operatorname{id}.
\]
$h^2=0$ は各項が互いに反交換し、それぞれの二乗が零であることによる。

明示的に偶奇へ折り畳むには
\[
 E_0=\bigoplus_{i\text{ even}}P_i(3i/2),\qquad
 E_1=\bigoplus_{i\text{ odd}}P_i(3(i-1)/2)
\]
とし、$d+h$ を
$D_1:E_1\to E_0$、$D_0:E_0(-3)\to E_1$ と読む。
上の内部 shift により両射は次数零であり、
$D_1D_0=F\operatorname{id}$、$D_0(3)D_1=F\operatorname{id}$ である。
偶数・奇数二項係数の和はそれぞれ $2^r$ なので階数の式を得る。

この明示 homotopy は Eisenbud の既知の Shamash 構成
\cite[Theorem~7.2, pp.~56--57]{NFMEisenbudCI}の仮定を満たす。
その自由分解の項は
\[
 \bigoplus_{j\ge0}P_{i-2j}\otimes_S(S/(F))(-3j)
\]
であり、十分高い次数の周期部分は、上の折畳みの同じ内部 shift を持つ。
従って graded singularity category でこの MF は $N$ の像を表す。
式 $D^2=+F$ が指定方程式に対する周期を固定しており、
方程式の unit を無言で掛け替えてはいない。
この周期は target 側で構成したもので、source 周期との比較 scalar は未計算である。
\end{proof}

\begin{lemma}[低 Hom の integral transfer を有限行列で検査する条件]
\label{lem:nfm-boundary-hom-flatness-test}
$C^\bullet$ を DVR $R$ 上の有界有限自由複体、$k$ を剰余体とする。
$b_i=\dim_KH^i(C_K)$ とし、$H^i(C)$ の非零巡回 torsion 直和因子の個数を $t_i$ とすると
\[
 \dim_kH^i(C\otimes_Rk)=b_i+t_i+t_{i+1}.
\]
従って全 $i$ で両 fibre の cohomology 次元が一致することは、
全 $H^i(C)$ が自由で底変更と可換になることに同値である。
このとき $C$ はその cohomology への $R$-線型 contraction を持つ。
有限個の orbit Hom 複体にこの条件を確認すれば、
それらの合成から必要な arity 二・三の homotopy transfer を $R$ 上で行える。
この判定に加えて三次 relation の原始性と表示比較の延長を確認しなければ、
周期位相の非分岐性は結論できない。
\end{lemma}

\begin{proof}
普遍係数の短完全列
\[
 0\to H^i(C)\otimes_Rk\to H^i(C\otimes_Rk)
    \to\operatorname{Tor}_1^R(H^{i+1}(C),k)\to0
\]
と DVR 上の有限生成加群の構造定理から次元式を得る。
全ての $t_i$ は非負なので次元の一致は torsion の消滅を意味する。
cohomology が自由なら、cycles と boundaries の二つの短完全列は
DVR 上の自由加群間の列として分裂する。
これらの分裂は cohomology の代表の包含、射影と contraction を与える。
積をこの contraction へ移す通常の有限 tree 式で $m_2,m_3$ を得る。
分母を導入しないのは、初めに $R$ 上の分裂を得たためである。
\end{proof}

\begin{corollary}[三次元境界の二次・三次積を integral に取れる]
\label{cor:nfm-boundary-threefold-integral-low-products}
$n=3$ の上の模型では、単一の enhanced orbit 模型から、
Adams 次数を保つ $R$ 上の strictly unital minimal $A_\infty$ 模型を選べる。
周期評価に必要な自由加群の階数は
\[
 (S_R^{1,-1})^*\simeq R^5,\qquad
 (S_R^{2,-2})^*\simeq R^{10},\qquad
 S_R^{2,-3}=Rt_p
\]
であり、$m_2,m_3$ を $R$ 上で取れる。
従って次に必要なデータは有限な二写像
\[
 m_2^*:R^{10}\longrightarrow (R^5)^{\otimes2},\qquad
 m_3^*(t_p^*)\in (R^5)^{\otimes3},\qquad t_p^*(t_p)=1
\]
となる。minimal 模型の正準性は主張しない。
$m_2^*$ の像が二次関係の部分束を定めるための一定階数条件、
および二次関係で割った三次 class の原始性は別に検査する必要がある。
\end{corollary}
\begin{proof}
命題\ref{prop:nfm-boundary-threefold-low-hom}と実際の MF 周期により、
全ての orbit 次数の Hom cohomology が $R$-自由である。
相対窓は graded section の切断、双対、評価の cone から構成されるので、
ここでの有限対象と射には同じ dg 模型を用いられる。
単に三角圏としての同値から積の適合を推測するのではない。
$Z$ の有限 graded MF 模型を一つ取り、その実際の Adams graded Hom 複体を使う。
各 homological/internal 次数の成分は有限自由 $R$-加群であり、合成は同じ模型上にある。
DVR 上で自由加群の部分加群は自由であり、cohomology も自由なので、
cycles と boundaries の短完全列を各 Adams 次数で分裂できる。
単位を代表として固定し、その直和補空間側で分裂を選べば、
strict unit を保つ contraction を得る。
三つの周期剰余について選んだ contraction を、同じ MF の厳密な周期で反復する。
この一つの模型への homotopy transfer の有限 tree 式が
$R$ 上の積を与え、内部次数も保つ。
$t_p$ は identity に対応する生成元として固定できる。
低次数の階数は前命題の表から従う。
これは積の integral な存在を示すもので、表示した二写像の値を評価した証明ではない。
\end{proof}

\begin{corollary}[三分点被覆では境界位相の三候補を区別できない]
\label{cor:nfm-hesse-boundary-period-candidates-blind}
この trait 上の $D$ の引戻しを $D_R=D(r(p,2+z))$ と書くと、
\[
 [D_R]=\left[\frac{p^3+8}{2z}\right]^{-1},
 \qquad L=K(\sqrt[3]{D_R}).
\]
従って $K^\times/K^{\times3}\to L^\times/L^{\times3}$ の核は
$\langle[D_R]\rangle$ であり、条件付き候補 $1,[D_R],[D_R]^2$ はすべて同じ自明類へ移る。
\end{corollary}
\begin{proof}
補題\ref{lem:nfm-period-hesse-eight-boundary-lines}の一次係数から
\[
 D_R\frac{p^3+8}{2z}
    =-\frac{(p^3+8)^3}{54}\bigl(1+O(z)\bigr)
\]
であり、右辺は $K$ の三乗である。これが二つの表示を与える。
巡回三次拡大で三乗根を持つ $b\in K^\times$ に対し、
その根の Galois character は三つの character の一つである。
$\sqrt[3]{D_R}$ の適当な冪で割ると根は $K$ に属するから、
$[b]\in\langle[D_R]\rangle$ である。
\end{proof}

\begin{remark}[三分点被覆と境界計算の使い分け]
三分点被覆に移って周期を計算するだけでは、元の trait での位相を判別できない。
一方、線束そのものの相対行列は元の trait 上にある。
以下ではこの行列を right window に正確に対応させ、二次・三次積の一次 jet を
元の trait 上で計算する。Hom の平坦性、二次関係の飽和性、回復した三次 relation の
原始性、表示同型の延長は別々に検証する。
\end{remark}

% !TEX root = ../nc-cubic.tex
\subsubsection{境界の低次積・二次関係の飽和と規格化周期}
\label{subsubsec:nfm-boundary-products-period}

本節では、元の trait 上で指定周期を保った低次積を計算する。
全次元の orbit Hom は平坦である一方、特殊 fibre の二次積は階数を三つ失う。
この三つの関係を飽和すると平坦な表示が回復し、中心三次の最初の項と
滑らかな potential 三次の同型延長から、完備局所周期位相を決定できる。
固定 MF・window・一般点生成性・飽和性の入力は引き続き前提とする。

\begin{lemma}[全次元の右窓と点線束の section module]
\label{lem:nfm-boundary-right-window-module-alln}
$R$ を本境界の DVR、$S=R[X,Y,Z,u_1,\ldots,u_r]$、
$A=S/(F_R)$、$r=n-1\ge2$ とし、既存の右射影で
$Z_R=j_R^!i_*\mathcal M[1]$ を定める。
$N_0=\operatorname{coker}(L:R[X,Y,Z](-1)^3\to R[X,Y,Z]^3)$ を
中心変数が零として作用する $A$-加群と見なす。
このとき window $i=1$ の enhanced functor により
\[
 \Psi_1Z_R\simeq N_0(-1)[1]\quad\text{in }D_{\mathrm{sg}}^{\mathrm{gr}}(A),
 \qquad\Psi_1=\pi_{\mathrm{sg}}R\Gamma_{\ge1}
\]
である。これは閉 fibre と一般 fibre の両方で、
内部 shift と方程式 $F_R$ の指定 MF 周期を保つ同定である。
\end{lemma}

\begin{proof}
適用する窓定理の環論的条件を確認する。
$S$ は regular DVR 上の多項式環、$F_R$ は正則な斉次三次である。
従って $A$ は noetherian かつ Gorenstein であり、graded duality により
$R\operatorname{Hom}_A(R,A)\simeq R(r)[1-(n+2)]$ となる。
$A_0=R$ は有限 global dimension を持ち、可換 noetherian graded 環なので
条件 $\chi$ を満たす。
Burke--Stevenson\cite[Definition~5.1, Proposition~4.18, Theorem~6.4(1)]{NFMBurkeStevenson}
を parameter $a=r$、window $i=1$ で適用すると
\[
 D^b(\operatorname{coh}\mathcal Y)
 =\langle\mathcal O(-r),\ldots,\mathcal O(-1),\mathcal B_1\rangle
\]
を得る。相対双対性で
$\mathcal B_1\cap\operatorname{Perf}(\mathcal Y)=\mathcal A_R$ である。
grade truncation、導来双対、自由 summand の除去という証明中の構成を
dg quotient で行い、$\Psi_1$ と窓の enhanced 同定を固定する。
ここでは単なる triangulated 同値から高次積の整合を推論しない。

$1\le j\le r$ と $t\ge1$ に対し $t-j\ge1-r$ である。
三次超曲面の完全列から中間 cohomology は零であり、最高次も
$H^n(\mathcal O(t-j))\simeq H^0(\mathcal O(j-t-r))^*=0$ である。
従って
\[
 R\Gamma_{\ge1}\mathcal O(-j)=A(-j)
\]
であり、この対象は特異圏で零となる。
右射影の有限個の mutation triangles を $\Psi_1$ で移すと
$\Psi_1(j_Rj_R^!i_*\mathcal M)=\Psi_1(i_*\mathcal M)$ となる。

$t\ge1$ では $H^1(\mathcal C,\mathcal M(t))=0$ であり、
線型表示の完全列から section module は
\[
 R\Gamma_{\ge1}(i_*\mathcal M)
 =\bigoplus_{t\ge1}H^0(\mathcal C,\mathcal M(t))=N_0(-1)
\]
となる。よって所要の式が従う。
これらの section modules と各有限 cone は $R$ 上で定義され、
必要な cohomology は自由なので閉／一般底変更と両立する。
これは任意の特異圏についての包括的な底変更定理ではない。

隣接する窓の grade truncation と自由 summand の除去を比較すると、
内部 shift $(-1)$ は右 mutation $R_{\mathcal O(-1)}(-1)$ に対応する。
MF 側では同じ曲率 $F_R$ に対する $(-3)=[-2]$ を用いる。
この実際の周期を窓へ移しているので、scalar の規格化変更はない。
共通の内部 shift $(-1)$ と cohomological shift $[1]$ は、
標準の符号付き dg 同型によって endomorphism 代数を変えない。
特に単位を保ち、周期の cohomological degree は二なので周期の符号も変わらない。
\end{proof}

\begin{lemma}[零曲率の中心 Koszul 因子の dg 分離]
\label{lem:nfm-boundary-koszul-dg-separation}
$T=k[u_1,\ldots,u_r]$、
$K_T=(T\otimes\Lambda(\epsilon_1,\ldots,\epsilon_r),d=\sum u_i\iota_i)$ とし、
$\deg\epsilon_i=(-1,1)$ とする。
ここで二次数は cohomological degree と Adams degree である。
縮約 $\iota_i$ を $e_i$ の像とする写像
\[
 \Lambda_k(e_1,\ldots,e_r)\longrightarrow\operatorname{End}_T(K_T),
 \qquad\deg e_i=(1,-1)
\]
は unital bigraded dg 代数の quasi-isomorphism である。
active 三次 $H$ の有限 graded MF $P$ に対し、tensor は
\[
 \operatorname{End}(P)\otimes\Lambda_k(e_1,\ldots,e_r)
 \longrightarrow\operatorname{End}(P\widehat\otimes K_T)
\]
という dg 代数の quasi-isomorphism を与え、実際の曲率 $H$ の MF 周期を厳密に保つ。
\end{lemma}

\begin{proof}
縮約は互いに反交換し、$d$ と超可換するので、表示した写像は dg 代数の射である。
augmentation $K_T\to k$ への postcomposition は
$\operatorname{End}_T(K_T)\to\operatorname{Hom}_T(K_T,k)$ という複体の射を与える。
$K_T$ は有限自由分解なので、この射は quasi-isomorphism である。
右辺の微分は零で、縮約の像はその外積基底を与える。
従って元の dg 代数の射も quasi-isomorphism である。
finite MF と有限 Koszul 複体の tensor の Hom は、符号付き tensor Hom に一致する。
$k$ は体なので tensor は quasi-isomorphism を保つ。MF 周期は $P$ の周期と
$K_T$ の恒等写像の tensor であり、scalar 変更はない。
\end{proof}

\begin{remark}[高次積を分離する際の意味]
active な厳密単位的 minimal $A_\infty$ 模型を、次数付き可換 dg 代数
$\Lambda(e_1,\ldots,e_r)$ で係数拡大すれば、
\[
 m_k(a_1\otimes\lambda_1,\ldots,a_k\otimes\lambda_k)
 =\pm m_k^{\mathrm{act}}(a_1,\ldots,a_k)
            \otimes\lambda_1\cdots\lambda_k
\]
という full 模型を取れる。
とくに $k\ge3$ で入力に純粋な中心生成元 $1\otimes e_i$ があれば積は零である。
任意の別の contraction で同じ係数公式が文字通り成立するという主張ではなく、
指定周期を保つ $A_\infty$ 同値の下でこの模型を選べるという意味である。
\end{remark}

\begin{proposition}[分離境界の active 二次・三次積は消える]
\label{prop:nfm-separated-boundary-active-products}
$L=L_0$、$Q_a=-b\operatorname{adj}L$ とし、
\[
 \mathsf D=\begin{pmatrix}0&L\\Q_a&0\end{pmatrix},\qquad
 \mathsf D^2=bXYZ I_6
\]
を active MF の微分とする。
その Adams graded endomorphism dg 代数を $\mathscr E_a$ と書く。
$E_{ij}$ を三行三列の行列単位とすると、
$H^{1,-1}(\mathscr E_a)$ は次の三つの cocycle の class を基底に持つ：
\[
 \begin{aligned}
 \phi_a&=\begin{pmatrix}0&-E_{21}\\bYE_{31}&0\end{pmatrix},\quad
 \phi_b=\begin{pmatrix}0&-E_{31}\\b(XE_{21}+2ZE_{31})&0\end{pmatrix},\\
 \phi_c&=\begin{pmatrix}0&-E_{22}\\b(XE_{31}+ZE_{33})&0\end{pmatrix}.
 \end{aligned}
\]
さらに
\[
 H^{2,-2}(\mathscr E_a)=k\{X\operatorname{id},Y\operatorname{id},Z\operatorname{id}\},
 \qquad H^{1,-2}(\mathscr E_a)=0,
 \qquad H^{2,-3}(\mathscr E_a)=k\operatorname{id}.
\]
ここで最後の identity は指定 MF 周期の class である。
この三次周期で正規化した minimal 模型において
\[
 m_2\big|_{(H^{1,-1})^{\otimes2}}=0,\qquad
 m_3\big|_{(H^{1,-1})^{\otimes3}}=0.
\]
後者の消滅も minimal 模型の gauge に依存しない。
\end{proposition}
\begin{proof}
微分は $\delta(f)=\mathsf Df-(-1)^{|f|}f\mathsf D$ である。
$(1,-1)$ 次数の cochain は、右上に定数行列 $C$、左下に一次行列を持ち、
計三十六係数を持つ。左下を $bD$ と書けば cocycle 条件は
\[
 LD-C\operatorname{adj}L=0,\qquad
 DL-\operatorname{adj}L\,C=0.
\]
この百八行三十六列の係数行列は階数三十三であり、
三十三次小行列の行列式は $-256p^2$ である。
表示した三つの独立な解があるので階数の上界も一致する。
この次数の boundaries は次数の理由で零である。

$(2,-2)$ 次数の cochain は二つの一次対角 block $(A,B)$ であり、
条件は $LB-AL=0$、$\operatorname{adj}L\,A-B\operatorname{adj}L=0$ である。
百四十四行五十四列の係数行列に四十二次 minor $-64p^3$ がある。
九次元の boundaries は定数行列 $C$ による $(LC,CL)$ で、
その独立性は九次 minor $-8$、または $L$ の一般点での可逆性から従う。
さらに $X\operatorname{id},Y\operatorname{id},Z\operatorname{id}$ が cocycle である。
これらの線型結合 $\ell\operatorname{id}$ が boundary なら、
$LC=\ell I_3$ から $\ell^3=-XYZ\det C$ となる。
三つの相異なる一次素因子を比較すると $\ell=0$、従って $C=0$ である。
これで kernel は十二次元、cohomology は表示した三次元となる。

$(1,-2)$ の cochain は左下の定数行列だけであり、その微分は $(LC,CL)$
なので cocycle は零である。
$(2,-3)$ は二つの定数対角 block である。
対応する八十一行十八列の係数行列は十七次 minor $-8p^2$ を持ち、
identity が kernel を張る。この次数の boundaries は零である。
これらの minor の添字と係数行列の生成式は、自己完結した検証器
\texttt{computations/verify\_boundary\_low\_products.py} に含めた。

次に積を計算する。九つの順序付き積のうち、唯一零でないものは
\[
 \phi_c\phi_b=\delta T,\qquad
 T=\begin{pmatrix}0&0\\-bE_{31}&0\end{pmatrix},
 \qquad \deg T=(1,-2)
\]
である。その他の八積は cochain として零である。
従って $m_2$ は零となる。
しかも全 $i\in\{a,b,c\}$ に対し
\[
 \phi_iT=T\phi_i=0
\]
が cochain として成り立つ。
$H^{1,-2}=0$ なのでこの boundary の homotopy は一意であり、
$h(\delta T)=T$ とする contraction を選べる。
この contraction を他の次数と strict unit に両立して延ばす。
arity 三の transfer 式の二つの項は上の二つの積なので、共に零である。
これで指定周期成分を含む $m_3$ の消滅を得る。

一次成分が恒等の gauge による変化では、
$f_2(H^{1,-1},H^{1,-1})\subset H^{1,-2}=0$ であり、
残る項も $m_2$ の上の制限が零なので消える。
従って三次消滅は gauge に依存しない。
\end{proof}

\begin{proposition}[split 境界模型の全 orbit Hom 平坦性と二次積の階数落ち]
\label{prop:nfm-boundary-all-orbit-flat-rank-drop}
$n\ge3$、$r=n-1$、$N=n+2$ とする。
pole 消去を厳密な不等式で行い、特殊方程式が $F_0=bXYZ$ となる模型を取る。
右窓の対象 $Z_R$ と、その実際の MF 周期を固定する。
一般 fibre については既存の固定 Tate--Koszul 入力を仮定する。
このとき全ての整数 $i,q$ に対し、
\[
 S_R^{i,q}=H^i\operatorname{RHom}(Z_R,\rho^qZ_R)
\]
は $R$-自由で、閉 fibre と一般 fibre への底変更に両立し、
その階数は $\binom N{3i+2q}$ である。範囲外の二項係数は零とする。
しかし二次積
\[
 \mu_R:(S_R^{1,-1})^{\otimes2}\longrightarrow S_R^{2,-2}
\]
の特殊 fibre の階数は $3r+\binom r2=\binom N2-3$ である。
その一般 fibre は全射なので、$\operatorname{coker}\mu_R$ は
正整数 $a_1,a_2,a_3$ に対する
$\bigoplus_{i=1}^3R/(z^{a_i})$ という torsion 加群である。
指数は後の命題\ref{prop:nfm-boundary-first-jet-products}で全て一と決定する。
\end{proposition}

\begin{proof}
補題\ref{lem:nfm-boundary-right-window-module-alln}により、$Z_R$ は点線束の section module $N_0(-1)[1]$ の
stabilization で表され、その同定は指定周期を保つ。
特殊 fibre の MF は active MF と $K_T$ の tensor なので、補題\ref{lem:nfm-boundary-koszul-dg-separation}が適用される。

active 部では、$q=0$ の cohomology は次数零と一に一つずつ、
$q=-1$ は次数一に三つ、$q=-2$ は次数二に三つである。
最初は $H^*(\mathcal O_C)$、二番目は $H^*(\mathcal O_C(-1))$、
三番目は右 mutation の triangle と、余因子 pairing による全射
$\mu:H^0(M^{-1}(1))\otimes H^0(M(1))\to H^0(\mathcal O_C(2))$
の双対の完全列から従う。これは $\mu^*:k^6\hookrightarrow k^9$ の cokernel が三次元という計算であり、他の cohomology は零である。
$\sigma_C^3=[-2]$ により、全 active orbit cohomology の二次数ベクトル空間は
$\Lambda(k^3)\otimes k[t,t^{-1}]$、$\deg t=(2,-3)$ と一致する。
これは積の同定ではない。
中心外積因子を tensor すると全特殊 fibre のベクトル空間は
$\Lambda(k^N)\otimes k[t,t^{-1}]$ となる。
一般 fibre は固定入力の $B^!\otimes K[t,t^{-1}]$ であり、
$H_{B^!}(s)=(1+s)^N$ なので同じ二次数次元を持つ。
各固定 $q$ の Hom は proper perfect な有限複体であるから、
DVR の普遍係数完全列により、その cohomology は自由となる。
$\deg t=(2,-3)$ から外積次数は $3i+2q$ であり、表示した階数式を得る。

特殊 fibre の次数一は active 三次元と中心 $r$ 次元に分かれ、
次数二は active 三次元、混合 $3r$ 次元、中心外積 $\binom r2$ 次元に分かれる。
active 二次積は命題\ref{prop:nfm-separated-boundary-active-products}により零で、残る二つの summand へは
tensor 積と外積が全射である。従って階数は $3r+\binom r2$ である。
一般 fibre の二次積は quadratic Koszul algebra の双対の次数二への積なので全射である。
従って cokernel は torsion であり、その剰余次元は三である。
DVR 上の構造定理から、非零な巡回 torsion 因子の個数は正確に三と分かる。
\end{proof}

\begin{corollary}[全次元の integral な低次積と原始的 target 周期]
\label{cor:nfm-boundary-alln-integral-products}
前命題の模型では、全 $n\ge3$ について、実際の MF 周期を固定した
strictly unital、Adams graded な $R$ 上の minimal $A_\infty$ 模型を取れる。
$N=n+2$ とすると
\[
 (S_R^{1,-1})^*\simeq R^N,\qquad
 (S_R^{2,-2})^*\simeq R^{\binom N2},\qquad
 S_R^{2,-3}=Rt_p
\]
であり、$t_p$ は原始的である。
\end{corollary}
\begin{proof}
全 cohomology の自由性が得られたので、
系\ref{cor:nfm-boundary-threefold-integral-low-products}の単一の enhanced MF 模型上の
contraction と finite tree の構成を全 $n$ に適用できる。
単位と周期を固定し、三つの Adams 周期剰余で分裂を選べばよい。
周期は $H^0\operatorname{End}(Z_R)=R\operatorname{id}$ から来るので原始的である。
ここでは target の周期と回復中心三次の原始性を同一視しない。
\end{proof}

\begin{lemma}[full 模型と active 計算の一次 jet の一致]
\label{lem:nfm-boundary-full-active-first-jet}
中心座標の pole 消去をさらに強めて
\[
 F_R\equiv J_R\pmod{z^2},\qquad F_R-J_R=\sum u_iG_i,\quad G_i\in z^2R[X,Y,Z,u]_2
\]
とできる。このとき、上の $R$ 上の minimal 模型を、$R/(z^2)$ 上で
active な相対行列 $L(q=2+z)$ の模型と $\Lambda(e_1,\ldots,e_r)$ との
係数拡大に一致するように取れる。この同定は実際の指定周期を保つ。
\end{lemma}
\begin{proof}
命題\ref{prop:nfm-hesse-boundary-integral-residual-model}の指数を、全係数について
$v_z(c)+Nj\ge2$ を満たすまで増やせばよい。一般 fibre の位相は変わらない。
補題\ref{lem:nfm-boundary-line-mf}の actual MF は、$R/(z^2)$ 上で active MF と
零曲率の中心 Koszul 複体の tensor に厳密に一致する。
補題\ref{lem:nfm-boundary-koszul-dg-separation}の外積代数からの写像は、
この非被約基底でも quasi-isomorphism である。実際、変数 $u_i$ は任意の基底上で
正則列であり、有限自由 Koszul 分解の augmentation に Hom を施す同じ証明が使える。
各内部次数でその cone は有界な自由複体で分裂するため、active Hom との tensor にも
quasi-isomorphism が保たれる。

全 cohomology が $R$-自由なので、各 Hom 複体の cycles と boundaries は自由で、
それらを定める短完全列は分裂する。この分裂の $R/(z^2)$ 上での任意の指定を、
自由基底を持ち上げて $R$ 上の分裂へ延ばせる。
この操作を同じ MF の三つの周期剰余で行い、単位を固定して周期的に反復する。
まず tensor dg 包含の像への retract を、その像上で homotopy が零となるように取り、
次に active contraction を外積代数で係数拡大する。
active 側の厳密単位性により、純粋な中心入力を含む三次積も零となる。
従って係数拡大による contraction と一致する一次 jet を持つ contraction を選べる。
有限 tree 式もこの底変更と可換なので、二次・三次積の一次 jet が一致する。
これは任意に選んだ別々の Hom 模型の積を貼り合わせる議論ではない。
\end{proof}

\begin{proposition}[消えた三つの二次関係は一次で回復する]
\label{prop:nfm-boundary-first-jet-products}
境界パラメータを $\varepsilon=q-2$ と書き、$d=p^3+8$、$b_0=d/3$ と置く。
補題\ref{lem:nfm-hesse-boundary-explicit-line}の実際の相対行列
$L(\varepsilon)$ と $-b_0\operatorname{adj}L(\varepsilon)$ を用いる。
命題\ref{prop:nfm-separated-boundary-active-products}の三つの active cocycle を $\phi_a,\phi_b,\phi_c$ と書き、
その定数上右 block の $(2,1),(3,1),(2,2)$ 成分を固定して一次まで持ち上げる。
$H^{2,-2}$ の基底を $(X\operatorname{id},Y\operatorname{id},Z\operatorname{id})$
とすると、列を $(aa,ab,ac,ba,bb,bc,ca,cb,cc)$ の順に並べた積は
\[
 m_2=\frac{\varepsilon}{3}
 \begin{pmatrix}
 0&0&p^{-1}&0&p&0&0&0&0\\
 p^{-1}&0&0&0&0&0&0&1&0\\
 0&0&0&2p^{-1}&0&0&0&0&p/4
 \end{pmatrix}+O(\varepsilon^2).
\]
従って三つの active elementary divisor の付値は全て一である。
飽和した二次 relation の特殊 fibre は
\[
 ac+p^2b^2,\qquad a^2+p\,cb,\qquad ba+\frac{p^2}{8}c^2
\]
である。ここで $a,b,c$ は cohomology の双対生成元である。
生成元を $x=a$、$y=pb$、$z=pc/2$ と改めると、これらは
\[
 yx=-\frac p2z^2,\qquad zy=-\frac p2x^2,\qquad xz=-\frac p2y^2
\]
という元の active 境界代数の関係に一致する。
\end{proposition}
\begin{proof}
$L=L_0+\varepsilon L_1+O(\varepsilon^2)$ とし、
$\phi_i=\begin{psmallmatrix}0&C_i\\b_0D_i&0\end{psmallmatrix}$ と書く。
$C_i=C_{i,0}+\varepsilon C_{i,1}+O(\varepsilon^2)$、
$D_i=D_{i,0}+\varepsilon D_{i,1}+O(\varepsilon^2)$ とする。
$J=\operatorname{adj}L$ と置けば、閉性は
\[
 LD_i-C_iJ=0,\qquad D_iL-JC_i=0
\]
であり、その一次式は特殊 fibre と同じ階数三十三の線型系である。
指定した三成分の正規化を加えると解が一意となる。
相対 $L$ の有理式を微分して解くと
\[
 \begin{aligned}
 C_{a,1}&=\frac1d\begin{pmatrix}-1&0&p^{-1}\\0&0&-2p^{-1}\\0&-2p^{-1}&0\end{pmatrix},
 \quad
 C_{b,1}=\frac1d\begin{pmatrix}p&p&0\\0&0&-2\\0&0&-2p^{-1}\end{pmatrix},\\
 C_{c,1}&=\frac1d\begin{pmatrix}-1&-1&0\\0&0&0\\0&-2p^{-1}&p/4\end{pmatrix}
 \end{aligned}
\]
を得る。また、積の計算に必要な下左 block の成分は
\[
 (D_{a,1})_{13}=-\frac{2Z}{pd},\qquad
 (D_{b,1})_{13}=-\frac{pX}{d},\qquad
 (D_{c,1})_{13}=-\frac{Y+2Z/p}{d}
\]
である。これらの式を含む全一次閉性は、有理係数の有限行列として検算できる。

次数 $(2,-2)$ の cocycle を $(A,B)$ とすると、その cohomology 座標は
\[
 ([X]A_{11},\ [Y]A_{11},\ [Z]A_{33})
\]
で与えられる。実際、boundary の上 block は $L_0H$ ($H$ 定数) なので、
$(L_0H)_{11}$ は $Z$ の倍数、$(L_0H)_{33}$ は $X,Y$ の線型結合であり、
この三成分は全 boundary を消す。また三つの scalar identity 上で標準基底となる。

唯一の非零な零次積 $\phi_c\phi_b=\delta T$ については、
$T$ の零次値を保ち、一次積から $\delta_1T$ を引いてからこの射影を用いる。
その他の積にはこの補正は不要である。
表示した $C_{i,1}$ と三つの $(D_{i,1})_{13}$ を代入すると上の行列となる。
三列 $(aa,ac,ba)$ の一次係数の determinant は
\[
 -\frac2{27p^3}\in k^\times
\]
である。従って、全係数が $\varepsilon$ で割れ、割った写像が全射なので、
Smith 付値はちょうど $(1,1,1)$ である。
双対の行を非零 scalar で規格化して、表示した三関係を得る。
最後の生成元変更は直接代入で確認できる。
\end{proof}

\begin{proposition}[飽和した二次関係が定める全次数の平坦族]
\label{prop:nfm-boundary-saturated-quadratic-flat}
上の積分模型で、中心を含む方程式係数が $z^2$ で割れるように
pole 消去の指数を取る。
$X_R=(S_R^{1,-1})^*$、$\mathcal R_2=\operatorname{im}(\mu_R^*)$、
\[
 \mathcal R_2^{\mathrm{sat}}
   =(\mathcal R_2\otimes_RK)\cap X_R^{\otimes2},
 \qquad
 B_R^{\mathrm{sat}}=T_R(X_R)/(\mathcal R_2^{\mathrm{sat}})
\]
と置く。
命題\ref{prop:nfm-boundary-first-jet-products}と full 一次 jet の一致により、三つの非単元 Smith 因子が全て $z$ であり、
閉 fibre の active 関係が
\[
 ac+p^2b^2=0,\qquad a^2+pcb=0,\qquad ba+(p^2/8)c^2=0
\]
となる。残りの関係は、補題\ref{lem:nfm-boundary-full-active-first-jet}により中心変数との可換関係である。
このとき $B_R^{\mathrm{sat}}$ の次数 $d$ の部分は有限自由で、階数は
$\binom{n+d+1}{d}$ である。
従って全次数で $R$-flat な quadratic algebra を得る。

また $B_R^{\mathrm{raw}}=T_R(X_R)/(\mathcal R_2)$ と置けば、
$B_R^{\mathrm{sat}}$ は $B_R^{\mathrm{raw}}$ をその $R$-torsion 全体で割ったものに等しい。
その torsion ideal は次数二の三つの元で生成され、$z$ で消える。
\end{proposition}

\begin{proof}
$\mathcal R_2^{\mathrm{sat}}$ は有限自由加群 $X_R^{\otimes2}$ 内の飽和部分加群なので、
その商は自由であり、部分加群の階数は一般 fibre と同じ $\binom{n+2}{2}$ である。
$X=a$、$Y=pb$、$Z=pc/2$ と置くと、active 関係は
\[
 ZY=-(p/2)X^2,\qquad XZ=-(p/2)Y^2,\qquad YX=-(p/2)Z^2
\]
となる。
$\upsilon=-p/2\in k^\times$ と置き、量子多項式環
\[
 T_{\upsilon}=k\langle X,Y,Z\rangle/
       (ZX-\upsilon XZ,\ XY-\upsilon YX,\ YZ-\upsilon ZY)
\]
と、その巡回自己同型 $\tau(X)=Z$、$\tau(Y)=X$、$\tau(Z)=Y$ を考える。
積 $a*b=a\tau^{\deg a}(b)$ による Zhang twist の二次関係は、ちょうど上の三式である。
これは tensor algebra の各次数で
$a_1\otimes\cdots\otimes a_d\mapsto
 a_1\tau(a_2)\cdots\tau^{d-1}(a_d)$ を用いても確認できる。
従って active 特殊代数の Hilbert 級数は $(1-t)^{-3}$ である。
中心変数 $r=n-1$ 個の多項式環を tensor するので、
全特殊代数の Hilbert 級数は $(1-t)^{-(n+2)}$ となる。

一方、一般 fibre の二次関係は固定 Tate--Koszul 入力による元の
quadratic AS-regular algebra $B_K$ の関係そのものである。
従って一般 fibre も Hilbert 級数 $(1-t)^{-(n+2)}$ を持つ。
各 $d$ で $(B_R^{\mathrm{sat}})_d$ は有限生成 $R$-加群である。
表示による代数の商は底変更に両立するので、その閉 fibre の次元も一般 rank も
$\binom{n+d+1}{d}$ である。
DVR 上の有限生成加群では、剰余次元は rank と非零巡回 torsion 因子数の和である。
従ってこの加群は torsion を持たず、有限自由である。
これが全次数について成立する。

$B_R^{\mathrm{raw}}\twoheadrightarrow B_R^{\mathrm{sat}}$ の核は
一般 fibre で零となるため $R$-torsion である。
逆に任意の torsion 元は torsion-free な標的へ零に写るので、核は torsion 全体である。
核は定義から $\mathcal R_2^{\mathrm{sat}}/\mathcal R_2$ の像で生成される ideal である。
計算した Smith 因子により、この次数二の商は $k^3$ で $z$ により消える。
$z$ は中心的なので、それが生成する全 ideal も $z$ で消える。
\end{proof}

\begin{proposition}[飽和二次族の中心商と原始的 potential]
\label{prop:nfm-boundary-saturated-active-potential}
命題\ref{prop:nfm-boundary-saturated-quadratic-flat}の平坦二次代数を $B_R$ と書く。
次数一の中心を $U_R=Z(B_R)_1$ とする。
このとき $U_R\subset(B_R)_1$ は自由階数 $r=n-1$ の直和因子であり、
その形成は任意の $R$-底変更に両立する。
その基底は中心正則列をなし、
\[
 \overline B_R=B_R/(U_R),\qquad V_R=(B_R)_1/U_R
\]
は各次数が有限自由な三生成元二次代数である。
一般 fibre は元の active 代数、閉 fibre は上の三角形 Sklyanin 代数となる。

さらに $\overline{\mathcal R}_2=\ker(V_R^{\otimes2}\to(\overline B_R)_2)$ と置くと、
\[
 \mathcal P_R=
 (\overline{\mathcal R}_2\otimes V_R)
        \cap(V_R\otimes\overline{\mathcal R}_2)
       \subset V_R^{\otimes3}
\]
は階数一の直和因子で、底変更に両立する。
元の active 代数の cyclic potential を一般 fibre の入力とすると、
$\mathcal P_R$ は cyclic invariant tensor からなる。
自然な交代商写像
$\operatorname{Alt}:V_R^{\otimes3}\to\bigwedge^3V_R$ は
$\mathcal P_R\xrightarrow{\sim}\bigwedge^3V_R$ を与える。
primitive な potential をこの写像で規格化すると、その symmetric 部分が定める
平面三次は $R$ 上で滑らかである。
これは点曲線 $V(J_R)$ ではなく、active 代数の potential に付随する三次である。
\end{proposition}

\begin{proof}
中心については有限自由加群間の写像
\[
 c:(B_R)_1\longrightarrow
     \operatorname{Hom}_R((B_R)_1,(B_R)_2),\qquad c(a)(b)=ab-ba
\]
を考える。一般 fibre の kernel は元の中心一次空間で次元 $r$ である。
閉 fibre の active 三角形代数には中心一次元がない。
実際、二次部分の基底を
$X^2,Y^2,Z^2,XY,YZ,ZX$ と取ると、
$[\alpha X+\beta Y+\gamma Z,X]$ の $XY,ZX$ 係数は
それぞれ $-\beta,\gamma$ であり、これが零なら $\beta=\gamma=0$ である。
さらに $Y$ と可換なら $\alpha=0$ となる。
従って閉 fibre の中心一次空間も、残りの多項式変数からなる次元 $r$ の空間である。
$c$ の一般・閉 fibre の rank は共に三であり、三次小行列式の一つは単元である。
一般 fibre で rank が三なので四次小行列式は $R$ 上でも全て零である。
単元小行列による消去から、kernel は直和因子で任意底変更に両立する。

$U_R$ の基底を選び、閉 fibre の中心変数の基底へ対応させる。
中心元 $u$ の閉 fibre が正則なら、$ua=0$ から $a\bmod z=0$ となる。
各次数の $R$-自由性を用いてこれを繰り返すと、
$a\in\bigcap_{j\ge0}z^j(B_R)_d=0$ である。
また $zb=ua$ なら、閉 fibre の正則性から $a=za_1$ と書け、
$z$ が正則なので $b=ua_1$ となる。
従って $B_R/(u)$ にも $R$-torsion はない。
各次数は有限なので、この商も次数ごとに有限自由である。
特殊代数では中心変数が正則列をなすため、この議論を順に適用できる。
これで中心正則列と商の全次数の自由性を得る。
構成した短完全列の商は $R$-flat なので、この正則列の計算も底変更と両立する。

$V_R$ の階数は三であり、二次商の階数は六である。
従って $\overline{\mathcal R}_2$ は自由階数三の直和因子である。
三次商は自由階数十なので、
$I_3=\ker(V_R^{\otimes3}\to(\overline B_R)_3)$ は自由階数十七の直和因子である。
二次 presentation から
\[
 (\overline{\mathcal R}_2\otimes V_R)\oplus
 (V_R\otimes\overline{\mathcal R}_2)\longrightarrow I_3,
 \qquad(a,b)\longmapsto a-b
\]
は全射で、その kernel は交わり $\mathcal P_R$ に同型である。
始域の階数は十八なので kernel は自由階数一で、全ての底変更に両立する。
その $V_R^{\otimes3}$ への包含も各 fibre で単射なので primitive である。
一般 fibre の cyclicity は、この tensor 空間の $R$-torsion-free 性により
$R$ 上へ延びる。

閉 fibre の関係を
$r_X=ZY+(p/2)X^2$、$r_Y=XZ+(p/2)Y^2$、
$r_Z=YX+(p/2)Z^2$ とする。
非零 tensor
\[
 w_0=Xr_X+Yr_Y+Zr_Z
     =r_XX+r_YY+r_ZZ
     =XZY+YXZ+ZYX+\frac p2(X^3+Y^3+Z^3)
\]
は交わりを生成し、
\[
 \operatorname{Alt}(w_0)=-3X\wedge Y\wedge Z,\qquad
 \operatorname{Sym}(w_0)
 =3XYZ+\frac p2(X^3+Y^3+Z^3)
 =\frac32 g_{p,2}
\]
である。ここで $\operatorname{Sym}$ も自然な symmetric 商写像とする。
$3$ は単元なので交代射は $R$ 上でも同型である。
例えば $\operatorname{Alt}(w_R)=X\wedge Y\wedge Z$ と規格化し、
$g_R=-2\operatorname{Sym}(w_R)$ と定義すれば、特殊三次は厳密に
\[
 g_R\bmod z=g_{p,2}=\frac p3(X^3+Y^3+Z^3)+2XYZ
\]
となる。
この三次が特異なら、勾配の共通非零零点で全座標が非零となり、
三つの勾配等式を掛けると $p^3=-8$ を得る。
本係数体では $p(p^3+8)$ は単元なので不可能である。
従って判別式は $R$ の単元となり、potential 三次は $R$ 上で滑らかである。
\end{proof}

\begin{proposition}[飽和二次関係に対する三次 class の最初の項]
\label{prop:nfm-boundary-first-cubic-class}
命題\ref{prop:nfm-boundary-first-jet-products}の相対模型において、指定 MF 周期を identity とし、
次数 $(2,-3)$ の cochain $\operatorname{diag}(A,B)$ からの射影を $A_{11}$ とする。
pair homotopy に上の正規化を用いた transfer では、
一次生成元上の三次テンソルの最初の項は
\[
 m_3^*(t^*)=\frac{\varepsilon}{3p}\,acb+O(\varepsilon^2).
\]
符号は補題\ref{lem:nfm-tate-cubic-relation-normalization}の
source 二拡大が $+\Omega$ を回復する規約に合わせた。
$\varepsilon$ で割った三次 class は、命題\ref{prop:nfm-boundary-first-jet-products}の飽和した二次関係による商で非零である。
従って、この飽和した模型における三次 relation の最初の消失次数は一である。
\end{proposition}
\begin{proof}
次数 $(2,-3)$ の boundaries は零、cycles は identity のみである。
従って $A_{11}=0$ という補空間を取れば、指定した射影は相対 contraction に延長できる。
積の homotopy の零次項は $(c,b)$ に対する
$T_{cb,0}=\begin{psmallmatrix}0&0\\-b_0E_{31}&0\end{psmallmatrix}$ だけである。
命題\ref{prop:nfm-separated-boundary-active-products}の
$\phi_iT=T\phi_i=0$ により、零次の三次 tree は全て零である。
従って $\Theta_0=0$ は商の class だけでなく tensor 自体として成立する。
また全ての $C_{i,0}$ の第一行は零である。
従って一次の $A_{11}$ 成分に寄与するのは
\[
 b_0\,(C_{i,1}(-E_{31}))_{11}
 \qquad\text{for }(j,k)=(c,b)
\]
だけである。命題\ref{prop:nfm-boundary-first-jet-products}の行列から、$i=a$ の場合だけ
$-b_0/(pd)=-1/(3p)$ となる。
この規約で寄与する tree の上 block は $-\phi_iT_{jk}$ なので、
主張の正符号を得る。規約の校正は一変数の
\[
 \mathsf D_x=\begin{pmatrix}0&x\\x^2&0\end{pmatrix},\qquad
 \phi=\begin{pmatrix}0&1\\-x&0\end{pmatrix},\qquad
 h=\begin{pmatrix}0&0\\-1&0\end{pmatrix}
\]
で行える。$\phi^2=\delta h$ であり、閉じた三次 tree
$-(h\phi+\phi h)=I_2$ は $+x^3$ を回復する。
この議論は同じ contraction の二つの有限 tree を計算しており、
独立に選んだ二次・三次積を混用していない。

三次 class の非零性は次の双対証明書でも確認できる。
三次語への線型汎関数を
\[
 \ell(aaa)=-p,\quad
 \ell(acb)=\ell(bac)=\ell(cba)=1,\quad
 \ell(bbb)=-p^{-2},\quad\ell(ccc)=-8p^{-2}
\]
とし、その他の語では零とする。
これは三つの二次関係の全ての左右一次倍を消し、$\ell(acb)=1$ である。
従って $[acb]\ne0$ である。
gauge による一次の変化は、一次二次関係の左右一次倍に入るので、
この商での非零性は gauge に依存しない。
\end{proof}

\begin{lemma}[滑らかな potential 三次による active 表示同型の延長]
\label{lem:nfm-boundary-potential-isom-extension}
$R$ を剰余体が標数零の DVR、$K$ をその商体とする。
$V,V'$ を階数三の自由 $R$-加群、
$Q\subset V^{\otimes2}$、$Q'\subset(V')^{\otimes2}$ を階数三の部分束とし、
$B=T_R(V)/(Q)$、$B'=T_R(V')/(Q')$ と置く。
$B_3,B'_3$ は自由階数十とする。
両一般 fibre の cubic syzygy line は cyclic tensor からなり、
両特殊 fibre は $S(0,-2,-p)$ と同型であると仮定する。
ここで剰余体において $p(p^3+8)\ne0$ とする。

このとき、任意の実際の次数付き $K$-同型
$f_K:B'_K\xrightarrow{\sim}B_K$ の次数一部分の射影類は
$R$ 上へ一意に延びる。その線型持上げ $f_R:V'\xrightarrow{\sim}V$ を選ぶと、
$f_R$ は両二次表示の $R$-同型であり、ある $\lambda\in K^\times$ に対して
$f_K|_{V'_K}=\lambda f_R|_{V'_K}$ となる。
\end{lemma}
\begin{proof}
まず
\[
 W=(Q\otimes V)\cap(V\otimes Q)\subset V^{\otimes3}
\]
を取る。
\[
 0\to W\to(Q\otimes V)\oplus(V\otimes Q)
   \to V^{\otimes3}\to B_3\to0
\]
と $B_3$ の自由性から、$W$ は自由階数一で、特殊 fibre と一般 fibre への
底変更に両立する。$Q$ が部分束なので二つの tensor 部分加群は飽和しており、
その交わり $W$ も飽和している。$W'$ も同様である。

$W_K$ 上の三巡回作用が恒等であるという仮定は、$V^{\otimes3}$ の
$R$-torsion-free 性により $W$ 上でも成立する。標数零での自然な分解
\[
 (V^{\otimes3})^{C_3}=\operatorname{Sym}^3V\oplus\bigwedge^3V
\]
を用いる。特殊 fibre $S(0,-2,-p)$ の cyclic potential は
\[
 w_0=-2(yxz+zyx+xzy)-p(x^3+y^3+z^3)
\]
である。その alternating 成分は
\[
 \varepsilon=xyz+yzx+zxy-xzy-zyx-yxz
\]
であって非零である。従って $W\to\bigwedge^3V$ は同型であり、
volume を固定すれば、alternating 成分が $\varepsilon$ となる生成元 $w$ を取れる。
同様に $w'$ を取る。

可換化写像 $\operatorname{ab}:V^{\otimes3}\to\operatorname{Sym}^3V$ により
$g=-\operatorname{ab}(w)/3$ と定めると、その特殊 fibre は適当な線型座標で
\[
 g_0=\frac p3(x^3+y^3+z^3)+2xyz
\]
の非零 scalar 倍となる。これは $p(p^3+8)\ne0$ により滑らかである。
従って $V(g)$ は $R$ 上の滑らかな平面三次であり、$g'$ も同様である。

$f_K$ は内在的に定まる $W'_K$ を $W_K$ へ移す。
alternating 成分を比較すると
\[
 f_K(w')=\det(f_K)w,\qquad f_K(g')=\det(f_K)g
\]
である。従って $f_K$ の射影類は、この二つの滑らかな偏極平面三次の同型である。
Benoist の Theorems~1.6(ii), 1.7 および Lemma~2.2(iii)
\cite[pp.~140--141]{NFMBenoistSeparation} により、
標数零の滑らかな平面三次の射影同型空間は有限であり、指定された一般点の
射影同型は $R$ 上へ一意に延びる。
ここで Theorem~1.7 の separatedness は対角の proper 性であり、
Theorem~1.6(ii) は $3$ を可逆にした基底上での Deligne--Mumford 性を与える。
三角形である点曲線の同型空間へ、この定理を適用しているのではない。

$\operatorname{Pic}(R)=0$ なので、得られた $\operatorname{PGL}_3(R)$ の点を
$f_R\in\operatorname{GL}_3(R)$ へ持ち上げられる。
$Q=Q_K\cap V^{\otimes2}$ と $Q'=Q'_K\cap(V')^{\otimes2}$ により、
一般 fibre での関係保存は $f_R(Q')=Q$ を与える。
従って $f_R$ は次数付き代数の同型である。
\end{proof}

\begin{lemma}[active 中心商での周期付値の評価]
\label{lem:nfm-boundary-active-lattice-valuation}
前補題の同型を用いる。$\Omega\in B_3$ は primitive とし、
$\Omega'\in B'_3$ は
\[
 f_K(\Omega')=a_K\Omega
\]
を満たすとする。$\Omega'$ の $B'_3$ における付値を $e$ とすると
\[
 [a_K]=[z^e u_R],\qquad u_R\in R^\times,
 \quad v_R([a_K])=e\pmod3.
\]
これは full 表示の同型の延長を要せず、実際の周期評価式を active 中心商へ
移せる場合に適用できる。
\end{lemma}
\begin{proof}
$f_K=\lambda f_R$ の次数三部分は $\lambda^3 f_R$ である。
$f_R$ は自由格子の同型なので、$f_R(\Omega')$ の付値も $e$ である。
$\Omega$ が primitive だから $f_R(\Omega')=z^e u_R\Omega$、$u_R\in R^\times$。
従って $a_K=\lambda^3z^e u_R$ である。
\end{proof}

\begin{lemma}[三角形 Sklyanin の自己同型と残余 scalar]
\label{lem:nfm-triangle-sklyanin-central-scalar}
$k$ は $\mu_3$ を含む標数零の体、$s=-p/2$、$p(p^3+8)\ne0$ とし、
\[
 S_0=k\langle X,Y,Z\rangle/(ZY-sX^2,\ XZ-sY^2,\ YX-sZ^2)
\]
と置く。この代数の射影次数付き自己同型群は $\mu_3^2$ であり、
任意の次数付き自己同型は $G=X^3=Y^3=Z^3$ を $k^{\times3}$ 倍する。
\end{lemma}
\begin{proof}
関係行列の determinant は $(1-s^3)XYZ$ であり、三角形上で階数は二である。
従って点対応は三成分を巡回する自己同型の graph である。
次数付き自己同型は三角形を保つので、その線型表示は monomial である。
成分置換は三巡回と可換でなければならず、従って巡回置換に限られる。
対角表示 $\operatorname{diag}(\alpha,\beta,\gamma)$ は
\[
 \beta\gamma=\alpha^2,\qquad
 \gamma\alpha=\beta^2,\qquad
 \alpha\beta=\gamma^2
\]
を満たす。射影的に $\alpha=1$ とすると $\beta^3=1$、$\gamma=\beta^2$ である。
従って射影群は、$\operatorname{diag}(1,\zeta,\zeta^2)$ と巡回置換による
位数九の群である。これらの線型持上げは $G$ を固定する。
残る grading scalar $\lambda$ は $G$ を $\lambda^3G$ へ移す。
\end{proof}

\begin{proposition}[全次元の完備局所周期位相]
\label{prop:nfm-boundary-period-local-kummer}
元の一ブロック一般族、正準半古典 frame、および固定 MF・window・
一般点生成性・coherent 飽和性の入力の下で、$n\ge3$ とする。
$k_0=\mathbb C(p,\mathbf m)$、$R=k_0[[z]]$、$K=k_0((z))$、$q=2+z$ と置く。
ここで $\mathbf m$ は任意個数の独立な混合・純中心係数である。
$D_R=D(r(p,2+z))$ とすると、正準強化点型 class の周期位相は
\begin{equation}\label{eq:nfm-boundary-period-local-kummer}
 \left.\gamma_{\mathrm{sc},n}\right|_K=[D_R]
   \quad\text{in }K^\times/K^{\times3},\qquad
 D_R=-\frac{(p^3+8)^2}{27}z+O(z^2).
\end{equation}
従ってこの torsor の局所体は $K(\sqrt[3]{D_R})$ であり、
$D=0$ の一般点での mod-three 付値は一、tame 分岐指数は三である。
この結論は全ての $n\ge3$ で同じである。
\end{proposition}

\begin{proof}
\emph{復元した三次の active 格子への移行。}
中心座標の pole 消去を $F_R\equiv J_R\pmod{z^2}$ となるまで行う。
これは一般 fibre 上の指定方程式と実際の周期を保つ座標変更であり、
体拡大や方程式生成元の自由な再規格化ではない。
命題\ref{prop:nfm-period-phase-low-products}により、一般 fibre の実際の
次数付き表示同型 $f_K$ に対して
\[
 \gamma_{\mathrm{sc},n}|_K=[a_K],\qquad
 f_K(\Omega_p)=a_K\Omega
\]
である。この式を次数一中心による商へ移す。
表示同型は内在的な中心を保ち、source 中心三次の active 像は非零なので、
同じ scalar $a_K$ を得る。

命題\ref{prop:nfm-boundary-saturated-quadratic-flat}と
\ref{prop:nfm-boundary-saturated-active-potential}により、target の飽和した
active 二次表示 $\overline B'_R$ は全次数で自由である。
ここで三次テンソルを $\Theta_R=m_3^*(t_p^*)$ と書く。
分離した零次模型では、active 三次 tree は文字通り零であり、
中心入力を含むものも厳密単位性により零である。
従って \emph{二次関係による商を取る前から}
$\Theta_R=z\omega_R$ と書ける。
中心商への射影を $\pi_R$ とすると
\[
 \left(\frac{\pi_R[\Theta_R]}z\right)\bmod z
      =\pi_0[\omega_R\bmod z].
\]
中心部分束の一次の傾きや、飽和関係の次の jet はこの式に寄与しない。
補題\ref{lem:nfm-boundary-full-active-first-jet}と
命題\ref{prop:nfm-boundary-first-cubic-class}から
\[
 \overline\Omega_p=z\omega_R,\qquad
 \omega_R\bmod z=\frac1{3p}[acb]\ne0.
\]
従って飽和 active 格子における $\overline\Omega_p$ の付値は正確に一である。

\emph{source 中心三次の原始性。}
source active 代数も同じ $R$ 上の Hesse 表示を持つ。
一般 fibre と特殊 fibre $S(0,-2,-p)$ の Hilbert 級数が共に
$(1-t)^{-3}$ なので、上と同じ次数別の議論で全次数が自由である。
その cyclic potential の可換部分は
$g_{p,q}=p(X^3+Y^3+Z^3)/3+qXYZ$ である。
量子化論文\cite[Hesse 型の普遍混合中心三次式の定理]{CubicQuantization}の
厳密な中心公式を $\hbar'=p$、$t=q/p$ で評価し、全体を $p$ 倍すると、
中心変数を零とした像は
\[
 \overline\Omega=
 \frac{pU}{3}(X^3+Y^3+Z^3)+V\operatorname{sym}(XYZ),
\]
\[
 U=1+\frac{q^2}{2}+\frac{p^3q}{6}-\frac{q^4}{48},\qquad
 V=\frac{(4-q^2)(12q+q^3-8p^3)}{48}.
\]
$q=2$ では $U=(p^3+8)/3$、$V=0$ であり、
$G=X^3=Y^3=Z^3\ne0$ を用いて
\begin{equation}\label{eq:nfm-boundary-source-center-primitive}
 \overline\Omega\bmod z=\frac{p(p^3+8)}3G
\end{equation}
となる。従って source 三次は primitive である。
ここでは固定入力の中心公式を用いており、その中心性全体の独立な再証明を
今回の一次 jet の検算で置き換えてはいない。

\emph{指定一般点同型の延長と単元の比較。}
両 active 表示の特殊 fibre は $S(0,-2,-p)$ であり、その potential 三次は滑らかである。
補題\ref{lem:nfm-boundary-potential-isom-extension}を実際の $f_K$ に適用し、
$f_K=\lambda f_R$ となる格子同型 $f_R$ を得る。
次数三では $\lambda^3$ が作用するので、Kummer 類には影響しない。
full 表示の中心変数方向の同型が $R$ 上へ延びることは必要ない。

特殊 fibre の基準同型を
$a\mapsto X$、$b\mapsto Y/p$、$c\mapsto2Z/p$ とする。
$XZY=-(p/2)X^3$ より $acb\mapsto-G/p$ であり、
\[
 f_0(\omega_R\bmod z)=-\frac1{3p^2}G.
\]
$f_R\bmod z$ とこの基準同型との差は source 特殊代数の次数付き自己同型である。
補題\ref{lem:nfm-triangle-sklyanin-central-scalar}により、
その $G$ 上の scalar は $k_0^{\times3}$ に属する。
従って補題\ref{lem:nfm-boundary-active-lattice-valuation}と
式\eqref{eq:nfm-boundary-source-center-primitive}から
\[
 [a_K]=[zu_R],\qquad
 u_R\bmod z=-\frac1{p^3(p^3+8)}
                  \quad\bmod k_0^{\times3}.
\]
一方、$D_R/z$ の剰余との比は
\[
 \frac{-1/[p^3(p^3+8)]}{-(p^3+8)^2/27}
       =\left(\frac3{p(p^3+8)}\right)^3.
\]
$R$ は完備、$3$ は可逆なので、$1+zR$ の任意の元は $1+zR$ 内で一意な三乗根を持つ。
従って $[a_K]=[D_R]$ である。
付値一から Kummer 拡大は次数三で全分岐し、剰余標数零なので tame である。
さらに補題\ref{lem:nfm-period-hesse-eight-boundary-lines}で
Hesse 写像の $q=2$ における分岐指数は一なので、元の $D$-付値も一 mod~三となる。
\end{proof}

\begin{corollary}[正準位相体の次数と大域式への残る条件]
\label{cor:nfm-boundary-phase-degree-three}
標準 split 一般係数体 $k_n$ 上で、全 $n\ge3$ の正準強化点型の周期位相は非自明であり、
その $\mu_3$-torsor は連結な三次体 $M/k_n$ である。
既存の scalar gerbe を $\alpha$ とすると、指定周期を標識した実現の最小次数は
\[
 3\operatorname{ind}(\alpha_M)\in\{3,9\}.
\]
$n=3,4$ では、この式は全 full class に共通である。

さらに $\overline\gamma_n$ の torsor が
$\operatorname{Spec}\mathbb C[s,t,1/(\Delta D)]$ 上で有限 étale に延びるという
追加条件の下では、大域的に $\overline\gamma_n=[D]$ となる。
この追加条件と $\operatorname{ind}(\alpha_M)$ の値は本計算では未決定である。
\end{corollary}
\begin{proof}
体拡大後に非自明なので、元の位相も非自明である。
$\mu_3\subset k_n$ かつ三が素数なので、対応する有限 étale 三次代数は体である。
標識付き最小次数の既存公式 $[M:k_n]\operatorname{ind}(\alpha_M)$ を適用する。
元の scalar index は三なので、拡大後の index は一または三である。
低次元での class 不変性は既存の周期整合定理による。
追加の非分岐条件の下では補題\ref{lem:nfm-period-two-boundary-criterion}により
候補が $1,[D],[D]^2$ に限られ、今回の $D$-付値一が $[D]$ を選ぶ。
\end{proof}

\begin{remark}[局所式の範囲と異なる二つの被覆]
正確な Kummer 等式\eqref{eq:nfm-boundary-period-local-kummer}は完備 trait 上の主張である。
未完備 trait 上でも mod-three 付値は一だが、剰余が三乗である単元が
その商体で三乗となることまでは主張しない。
また命題\ref{prop:nfm-hesse-boundary-third-division}の三分点 torsor は、
この局所体の三つのコピーからなる九次被覆である。
正準位相の三次被覆とは、一成分の体が一致するだけであり、torsor 全体は異なる。
全次元での局所式の一致から、二変数不変量体上での位相の一致や
他の divisor での非分岐性を結論してはいない。
任意の方程式 frame の変更則や独立標的尺度による次数九の例も引き続き有効である。
\end{remark}

